\documentclass[11pt,reqno]{amsart}
\usepackage[localtheorems]{raeez-math-template}

% ================================================================
% Packages
% ================================================================
\usepackage{array}

% ================================================================
% Theorem environments
% ================================================================
\newtheorem{theorem}{Theorem}[section]
\newtheorem{proposition}[theorem]{Proposition}
\newtheorem{lemma}[theorem]{Lemma}
\newtheorem{corollary}[theorem]{Corollary}
\newtheorem{conjecture}[theorem]{Conjecture}
\theoremstyle{definition}
\newtheorem{definition}[theorem]{Definition}
\newtheorem{construction}[theorem]{Construction}
\newtheorem{example}[theorem]{Example}
\newtheorem{warning}[theorem]{Warning}
\theoremstyle{remark}
\newtheorem{remark}[theorem]{Remark}
\newtheorem{principle}[theorem]{Principle}

% ================================================================
% Macros
% ================================================================
\newcommand{\cA}{\mathcal{A}}
\newcommand{\cB}{\mathcal{B}}
\newcommand{\cC}{\mathcal{C}}
\newcommand{\cD}{\mathcal{D}}
\newcommand{\cF}{\mathcal{F}}
\newcommand{\cH}{\mathcal{H}}
\newcommand{\cM}{\mathcal{M}}
\newcommand{\cO}{\mathcal{O}}
\newcommand{\cL}{\mathcal{L}}
\newcommand{\cW}{\mathcal{W}}
\newcommand{\cP}{\mathcal{P}}
\newcommand{\cZ}{\mathcal{Z}}
\newcommand{\cR}{\mathcal{R}}
\newcommand{\SCchtop}{\mathsf{SC}^{\mathrm{ch,top}}}

\newcommand{\Ran}{\mathrm{Ran}}
\newcommand{\MC}{\mathrm{MC}}
\newcommand{\HH}{\mathrm{HH}}
\newcommand{\CE}{\mathrm{CE}}
\newcommand{\ChirHoch}{\mathrm{HH}^{*}_{\mathrm{ch}}}

\newcommand{\Ainf}{A_{\infty}}
\newcommand{\Eone}{\mathsf{E}_{1}}
\newcommand{\Etwo}{\mathsf{E}_{2}}
\newcommand{\Ethree}{\mathsf{E}_{3}}
\newcommand{\En}{\mathsf{E}_{n}}
\newcommand{\Einf}{\mathsf{E}_{\infty}}

\newcommand{\Sym}{\mathrm{Sym}}
\newcommand{\Ass}{\mathrm{Ass}}
\newcommand{\Com}{\mathrm{Com}}
\newcommand{\Lie}{\mathrm{Lie}}

\DeclareMathOperator{\Hom}{Hom}
\DeclareMathOperator{\RHom}{RHom}
\DeclareMathOperator{\Res}{Res}
\DeclareMathOperator{\Aut}{Aut}
\DeclareMathOperator{\End}{End}
\DeclareMathOperator{\Conf}{Conf}
\DeclareMathOperator{\FM}{FM}
\DeclareMathOperator{\Spec}{Spec}
\DeclareMathOperator{\av}{av}
\DeclareMathOperator{\rk}{rk}
\DeclareMathOperator{\GRT}{GRT}
\DeclareMathOperator{\PBr}{PBr}
\DeclareMathOperator{\depth}{depth}

\newcommand{\fg}{\mathfrak{g}}
\newcommand{\fh}{\mathfrak{h}}

\newcommand{\gmod}{\mathfrak{g}^{\mathrm{mod}}_{\cA}}
\newcommand{\gEone}{\mathfrak{g}^{\Eone}_{\cA}}

\newcommand{\Barord}{\overline{B}^{\mathrm{ord}}}
\newcommand{\BarSig}{\overline{B}^{\Sigma}}
\newcommand{\Barch}{\overline{B}^{\mathrm{ch}}}
\newcommand{\Cobar}{\Omega^{\mathrm{ch}}}

\newcommand{\Vir}{\mathrm{Vir}}
\newcommand{\KM}{\mathrm{KM}}
\newcommand{\BP}{\mathrm{BP}}

\newcommand{\CC}{\mathbb{C}}
\newcommand{\RR}{\mathbb{R}}
\newcommand{\ZZ}{\mathbb{Z}}
\newcommand{\PP}{\mathbb{P}}
\newcommand{\id}{\mathrm{id}}
\newcommand{\op}{\mathrm{op}}
\newcommand{\FMAss}{F\!\Ass}

\numberwithin{equation}{section}

% ================================================================
\begin{document}

\title[The Ordered Bar Complex and $\Eone$ Primacy]
{The Ordered Bar Complex and $\Eone$ Primacy}

\author{Raeez Lorgat}
\address{Perimeter Institute for Theoretical Physics,
  Waterloo, ON N2L 2Y5, Canada \\
  Department of Physics, University of Waterloo,
  Waterloo, ON N2L 3G1, Canada}
\email{rlorgat@perimeterinstitute.ca}

\date{}

\begin{abstract}
The five theorems of modular Koszul duality (bar-cobar adjunction,
inversion, complementarity, genus universality, chiral Hochschild
concentration) are proved in the literature for the symmetric bar
complex $\BarSig(\cA)$. The symmetric bar is a shadow: the
$\Sigma_n$-coinvariant projection of the ordered bar complex
$\Barord(\cA) = T^c(s^{-1}\bar\cA)$, whose deconcatenation
coproduct carries the full $\Eone$-chiral coalgebra structure.
This paper establishes $\Eone$ primacy: the ordered bar is the
primitive object; the symmetric bar is derived by averaging; and the
five theorems are the invariants that survive the quotient.

After a strong-unitary \(R\)-twisted descent datum is fixed, the
averaging map $\av_R\colon\gEone\twoheadrightarrow\gmod$ is
a surjective morphism of dg~Lie algebras that sends the ordered
Maurer--Cartan element $\Theta^{\Eone}_\cA$ to the modular
$\Theta_\cA$, projects the classical $r$-matrix $r(z)$ to
$\kappa(\cA)$, and projects the Drinfeld--Kontsevich associator
$\Phi_{\mathrm{KZ}}$ to the cubic shadow~$\mathfrak{C}(\cA)$.
The map does not split as dg~Lie algebras: the kernel at degree~$3$
contains the linearised associator $[\Omega_{12},\Omega_{23}]$,
and the space of Lie sections compatible with the Maurer--Cartan
equation is a torsor for the Grothendieck--Teichm\"uller group
$\GRT_1$. At genus~$1$, a fourth obstruction to splitting arises
from the quasi-modular anomaly $E_2 \mapsto \widehat{E}_2$.

For the principal examples: Heisenberg $\cH_k$ has
$r(z) = k/z$ and $\kappa(\cH_k) = k$; affine Kac--Moody
$\widehat{\fg}_k$ has $r(z) = k\Omega_{\fg}/z$ (trace-form
convention; vanishing at $k=0$) and
$\kappa(\widehat{\fg}_k) = \dim(\fg)(k+h^\vee)/(2h^\vee)$,
the Sugawara shift $\dim(\fg)/2$ arising from the simple-pole
self-contraction; Virasoro $\Vir_c$ has
$r(z) = (c/2)/z^3 + 2T/z$ and $\kappa(\Vir_c) = c/2$.
The non-trivial non-coinvariant component of the KZ associator,
represented infinitesimally by $[\Omega_{12},\Omega_{23}]$, lies in
the kernel of \(\av_R\). Its trivial isotypic projection is the cubic
shadow \(\mathfrak C(\cA)\); this scalar shadow does not reconstruct
\(\Phi_{\mathrm{KZ}}\).

The bar direction does not by itself uplift a $2$-real-dimensional
chiral algebra to a $3$-real-dimensional holomorphic--topological
theory. The dimensional uplift is governed by the chiral
Deligne--Tamarkin theorem: the universal local Swiss-cheese pair
$(Z^{\mathrm{der}}_{\mathrm{ch}}(A_b),\, A_b)$ is terminal in the
category of two-coloured $\SCchtop$-pairs over~$A_b$, the open
colour~$A_b$ is acted on by the closed colour
$Z^{\mathrm{der}}_{\mathrm{ch}}(A_b)$ through the chiral
Swiss-cheese mixed operations, and the closed sector is not acted
on by the open sector. The bar coalgebra
$\Barord(\cA) = T^c(s^{-1}\bar\cA)$ is a chain-level
\emph{computational model} that resolves $\cA$ for the purpose of
computing $Z^{\mathrm{der}}_{\mathrm{ch}}(\cA)
= \HH^\bullet_{\mathrm{ch}}(\cA, \cA)
= \mathrm{Hom}_{\mathrm{ch}}(\Barord(\cA), \cA)$; it is not the
source of the dimensional uplift, which is the universal property
of the Swiss-cheese pair. We prove the $\Eone$ primacy theorem,
state the five ordered counterparts
$\mathrm{A}^{\Eone}$--$\mathrm{H}^{\Eone}$, exhibit the chiral
Deligne--Tamarkin Swiss-cheese theorem with its hypothesis package,
and exhibit the formality bridge and its failure.
\end{abstract}

\subjclass[2020]{17B69, 18M70, 18M85, 81T40}

\keywords{Ordered bar complex, $\Eone$-algebras, chiral algebras,
Koszul duality, $R$-matrix, Drinfeld associator, modular operad,
averaging map, Grothendieck--Teichm\"uller group, Yangian}

\maketitle

% ================================================================
% SECTION 1
% ================================================================
\section{Introduction: the averaging map and what it loses}
\label{sec:introduction}

The bar complex attached to an augmented associative algebra admits
two very different presentations. In the first, one takes the reduced
tensor coalgebra $T^c(s^{-1}\bar A)$ with its deconcatenation
coproduct; the bar differential is built from the associative
product in consecutive slots. In the second, one passes to
$\Sigma_n$-coinvariants and works with the symmetric bar complex,
which is the natural home for commutative or $\Einf$-algebras and
for the Francis--Gaitsgory chiral bar. The historical order in
chiral algebra and factorization homology has been: symmetric bar
first, ordered refinement second. This paper reverses the order of
presentation, and with it the logical role of the two objects.

\subsection{The thesis}

The thesis is that the ordered bar complex
$\Barord(\cA) = T^c(s^{-1}\bar\cA)$ is the primitive object of
modular Koszul duality, and that the symmetric bar
$\BarSig(\cA)$ is an $R$-twisted coinvariant shadow when descent
exists; in the pole-free locus this reduces to ordinary averaging. The
map $\av_R$ is a surjective morphism of dg~Lie algebras after a
strong-unitary \(R\)-twisted descent datum is fixed. Under the same
descent hypotheses, the five foundational theorems of modular Koszul
duality (Theorems A through D, together with the chiral Hochschild
theorem~H) have ordered counterparts whose symmetric forms are their
\(R\)-coinvariant images. The information lost under $\av_R$ is geometric,
arithmetic, and physical: the spectral $r$-matrix, the KZ connection,
the Drinfeld associator, and the full braiding of tensor categories
of modules all live on the ordered side.

Three principles motivate the inversion. The first is operadic:
Stasheff's associahedra $K_n$ \cite{Stasheff63} carry more
information than the symmetric groups $\Sigma_n$, and the
Getzler--Jones \cite{GJ94} theory of $\En$-operads makes $\Eone$ the
lowest rung of a ladder whose every step is non-commutative. The
second is representation-theoretic: the classical $r$-matrix of an
affine Kac--Moody algebra at level~$k$ is the content of the
degree-$2$ genus-$0$ Maurer--Cartan element on the ordered side, and
its $\Sigma_2$-average recovers the degree-$2$ scalar shadow: for
Heisenberg this is~$\kappa$, while for non-abelian affine Kac--Moody
it is $\kappa_{\mathrm{dp}} = k\dim(\fg)/(2h^\vee)$, with the full
$\kappa$ obtained after adding the Sugawara shift $\dim(\fg)/2$. The
third is operadic--holographic: the chiral Deligne--Tamarkin
theorem (Theorem~\ref{thm:swiss-cheese-chiral-DT-standalone}
below) constructs the universal Swiss-cheese pair
$(Z^{\mathrm{der}}_{\mathrm{ch}}(\cA),\, \cA)$ as terminal in the
category of two-coloured $\SCchtop$-pairs over $\cA$. The closed
colour $Z^{\mathrm{der}}_{\mathrm{ch}}(\cA)$ acts on the open
colour $\cA$ through the chiral Swiss-cheese mixed operations; the
open colour does not act back. The bar coalgebra
$\Barord(\cA)$ is a chain-level computational model that resolves
$\cA$ for the purpose of computing
$Z^{\mathrm{der}}_{\mathrm{ch}}(\cA)$ via
$\mathrm{Hom}_{\mathrm{ch}}(\Barord(\cA), \cA)$; the $2$-real to
$3$-real holomorphic--topological dimensional uplift is a property
of the Swiss-cheese pair and the universal collar of its open
boundary, not of the bar functor applied to $\cA$ alone. The
averaging map sends the open ordered structure to its symmetric
$\Sigma_n$-coinvariant image; the closed--open Swiss-cheese
relation is invisible to averaging on a single colour and requires
the two-coloured operad.

We work throughout with the conventions of the monograph~\cite{Lor26},
which we refer to as Volume~I. Everything is stated with OPE modes
rather than lambda-brackets; conformal weights are denoted $h$, and
desuspended degrees $|s^{-1}v| = |v| - 1$ (the bar complex lowers
degree). The grading is cohomological.

\subsection{Notation and conventions}
\label{subsec:notation}

Let $k$ be a field of characteristic zero, $X$ a smooth algebraic
curve over~$k$, and $\cA$ an augmented cyclic chiral algebra on~$X$
in the sense of Beilinson--Drinfeld~\cite{BD04}. We write
$\bar\cA$ for the augmentation ideal
$\ker(\varepsilon\colon \cA \to k)$. All bar constructions are
defined on~$\bar\cA$, not on $\cA$ itself~\cite{Lor26}.
Symmetric groups $\Sigma_n$ act on tensor powers in the standard
way; coinvariants are left exact and are denoted with a lower
subscript, $(-)_{\Sigma_n}$. The ordered and symmetric
Fulton--MacPherson compactifications are written
$\overline{\FM}_n^{\mathrm{ord}}(\CC)$ and
$\overline{\FM}_n(\CC)$ respectively; the former is the
Ran-space blow-up along ordered collision diagonals.

% FOUR OBJECTS:
% 1. B(A) = bar coalgebra = T^c(s^{-1} A-bar) with deconcatenation + twist
% 2. A^i = H^*(B(A)) = dual coalgebra (Koszul cohomology of bar)
% 3. A^! = ((A^i)^v) = dual ALGEBRA (linear or Verdier dual)
% 4. Z^der_ch(A) = derived chiral center = Hochschild cochains = bulk

\subsection{Summary of results}

The main results are:

\begin{enumerate}[label=(\roman*)]
\item The $\Eone$ \emph{convolution dg~Lie algebra}
  $\gEone$ (Definition~\ref{def:e1-convolution}) and its
  ordered Maurer--Cartan element $\Theta^{\Eone}_\cA$
  (Theorem~\ref{thm:e1-mc}).

\item The \emph{averaging map} $\av_R\colon \gEone
  \twoheadrightarrow \gmod$ as a surjective dg~Lie morphism
  (Theorem~\ref{thm:av-surjective}), with the projection
  formula $\av_R(\Theta^{\Eone}_\cA) = \Theta_\cA$
  (Theorem~\ref{thm:av-mc-projection}).

\item The \emph{$\Eone$ primacy theorem}
  (Theorem~\ref{thm:e1-primacy-main}): surjectivity, MC
  projection, bracket preservation, and non-splitting of $\av_R$.

\item The \emph{non-splitting proposition}
  (Proposition~\ref{prop:nonsplitting}): the kernel of $\av_R$
  is controlled by $\GRT_1$ at genus~$0$ and by the
  quasi-modular anomaly at genus~$1$.

\item The \emph{five ordered theorems}
  $\mathrm{A}^{\Eone}$--$\mathrm{H}^{\Eone}$
  (Theorems~\ref{thm:ordered-A}--\ref{thm:ordered-H}),
  whose $\av_R$-images are the classical five theorems when
  \(R\)-twisted descent is fixed.

\item Complete degree-$2$ computations for Heisenberg, affine
  Kac--Moody, and Virasoro (Section~\ref{sec:heis-km}).
\end{enumerate}

% ================================================================
% SECTION 2
% ================================================================
\section{The ordered bar complex}
\label{sec:ordered-bar}

\subsection{Three coalgebra structures on $T(V)$}

Before defining the ordered bar complex we distinguish three
cooperadic structures carried by a sum of tensor powers. All three
arise somewhere in modular Koszul duality; only one equips the
ordered bar.

\begin{definition}[The three coalgebra structures]
\label{def:three-coalgebras}
Let $V$ be a graded vector space and
$T(V) = \bigoplus_{n \geq 0} V^{\otimes n}$ its tensor space.
Three cooperad structures arise naturally:
\begin{enumerate}[label=(\roman*)]
\item The \emph{Harrison coLie} cooperad $\Lie^c$, whose
  degree-$n$ component has dimension $(n-1)!$ and whose
  coproduct decomposes a Lie word into pairs of sub-Lie-words
  using the Harrison differential.

\item The \emph{coshuffle cooperad} $\Sym^c$, whose degree-$n$
  coproduct produces $2^n$ summands by the shuffle formula
  \begin{equation}\label{eq:coshuffle}
    \Delta_{\mathrm{sh}}(v_1 \cdots v_n)
    = \sum_{I \sqcup J = \{1,\ldots,n\}}
      \varepsilon(I,J)\, (v_I) \otimes (v_J),
  \end{equation}
  with the sum running over all $2^n$ bipartitions and
  $\varepsilon(I,J)$ the Koszul sign.

\item The \emph{deconcatenation cooperad} $T^c$, whose
  degree-$n$ coproduct produces $n+1$ summands by splitting
  the word at a single position:
  \begin{equation}\label{eq:deconc}
    \Delta(v_1 \cdots v_n)
    = \sum_{i=0}^{n}
      (v_1 \cdots v_i) \otimes (v_{i+1} \cdots v_n).
  \end{equation}
\end{enumerate}
\end{definition}

The three structures are \emph{not} interconvertible without loss
of information. At degree~$3$ the Harrison differential sees $2$
terms, the coshuffle sees $8$, and deconcatenation sees~$4$. Only
the last is ordered and strictly associative at the coalgebra level.

\begin{proposition}[Coshuffle is not deconcatenation]
\label{prop:coshuffle-neq-deconc}
Let $V$ be nonzero. The coshuffle coproduct on $\Sym(V)$ and the
deconcatenation coproduct on $T^c(V)$ are distinct cooperadic
structures: they have different degree-$n$ term counts ($2^n$ vs.\
$n+1$), different symmetry properties (the coshuffle is
cocommutative; deconcatenation is not), and different primitive
subspaces.
\end{proposition}

\begin{proof}
Cocommutativity of the coshuffle follows from the involution
$I \leftrightarrow J$ on bipartitions. Deconcatenation does not
admit such an involution: swapping $(v_1 \cdots v_i)$ with
$(v_{i+1} \cdots v_n)$ does not preserve the ordered word.
The term counts are elementary.
\end{proof}

\begin{remark}[Where each structure appears]
\label{rem:where-coalgebras}
The Harrison coLie cooperad $\Lie^c$ underlies the chiral Lie bar
complex used by Francis--Gaitsgory~\cite{FG12}, where only the
zeroth product $a_{(0)}b$ of the OPE is visible. The coshuffle
cooperad $\Sym^c$ underlies the symmetric bar $\BarSig(\cA)$ on
the Ran space, where factorization is implemented by the topological
coproduct on unordered configuration space. The deconcatenation
cooperad $T^c$ underlies the ordered bar $\Barord(\cA)$, where
the coproduct splits a word along a single internal edge and
remembers the ordering on the collision points. The averaging map
$\av$ of Section~\ref{sec:averaging} descends from $T^c$ to
$\Sym^c$ by external $R$-twisted $\Sigma_n$-coinvariants
$\BarSig_n(\cA) \simeq (\Barord_n(\cA))^{R\text{-}\Sigma_n\text{-coinv}}$
(the $R$-twist is trivial on the pole-free subclass, where naive
$\Sigma_n$-coinvariants suffice;); it does not factor
through Harrison unless one further restricts to the Eulerian
weight-$1$ summand.
\end{remark}

\subsection{The ordered bar complex $\Barord(\cA)$}

\begin{definition}[The ordered bar complex]
\label{def:ordered-bar}
Let $\cA$ be a strongly admissible augmented $\Eone$-chiral algebra
on~$X$, with augmentation ideal $\bar\cA = \ker \varepsilon$. The
\emph{ordered bar complex} is the graded vector space
\begin{equation}\label{eq:bar-definition}
  \Barord(\cA) \;:=\; T^c(s^{-1}\bar\cA)
  \;=\; \bigoplus_{n \geq 0} (s^{-1}\bar\cA)^{\otimes n},
\end{equation}
equipped with:
\begin{enumerate}[label=(\alph*)]
\item The \emph{deconcatenation coproduct}
  \begin{equation}\label{eq:bar-coproduct}
    \Delta\bigl([a_1|\cdots|a_n]\bigr)
    \;=\; \sum_{i=0}^{n}
    [a_1|\cdots|a_i] \otimes [a_{i+1}|\cdots|a_n],
  \end{equation}
  making $\Barord(\cA)$ a cofree conilpotent coalgebra
  over~$T^c$.

\item The \emph{bar differential} $d_{\mathrm{bar}}$, obtained
  by extracting OPE collision residues at consecutive pairs of
  points in the ordered configuration space
  $\overline{\FM}_n^{\mathrm{ord}}(\CC)$:
  \begin{equation}\label{eq:bar-differential}
    d_{\mathrm{bar}}\bigl([a_1|\cdots|a_n]\bigr)
    \;=\; \sum_{i=1}^{n-1} \pm\,
    [a_1|\cdots|a_i \star a_{i+1}|\cdots|a_n],
  \end{equation}
  where $a_i \star a_{i+1}$ denotes the collision residue
  (the coefficient of the simple pole in the OPE of $a_i(z)$
  and $a_{i+1}(w)$ against the Arnold form
  $\eta = d\log(z-w)$).
\end{enumerate}
The desuspension $s^{-1}$ lowers the degree by one:
$|s^{-1}a| = |a| - 1$. The augmentation ideal is essential;
using $\cA$ in place of $\bar\cA$ would include the unit and
break the construction.
\end{definition}

\begin{theorem}[Ordered bar differential squared]
\label{thm:bar-d-squared}
The ordered bar differential satisfies $d_{\mathrm{bar}}^2 = 0$.
\end{theorem}

\begin{proof}
The argument is a codimension-two face pairing on
$\overline{\FM}_n^{\mathrm{ord}}(\CC)$ with the Arnold relation
supplying the required cancellation between the two ways of
contracting an adjacent pair of collisions. Each codimension-$2$
stratum is reached by two codimension-$1$ paths (contract $e_1$ then
$e_2$, or $e_2$ then $e_1$), and the edge-orientation determinant
supplies a minus sign between them. The cancellation uses only the
associativity of face composition in the face poset
(Getzler--Kapranov~\cite{GK98}), which holds for ordered
configurations by the same argument as for unordered stable graphs.
\end{proof}

\subsection{The classical $r$-matrix as degree-$2$ bar data}

The binary collision residue of the ordered bar differential extracts
the classical $r$-matrix. This is the degree-$2$ component of the
$\Eone$ Maurer--Cartan element, and it is the central object of the
$\Eone$ theory.

\begin{definition}[The classical $r$-matrix from the bar complex]
\label{def:r-matrix-from-bar}
Let $\cA$ be a cyclic chiral algebra on~$X$ with OPE
$a(z)\,b(w) \sim \sum_{n \geq 0} (a_{(n)}b)(w)/(z-w)^{n+1}$.
The \emph{classical $r$-matrix} of $\cA$ is the degree-$2$
genus-$0$ component of the $\Eone$ Maurer--Cartan element:
\begin{equation}\label{eq:r-matrix-definition}
  r_\cA(z) \;:=\;
  \Res^{\mathrm{coll}}_{0,2}(\Theta^{\Eone}_\cA)
  \;\in\; \End_\cA(2) \otimes \cO(*\Delta).
\end{equation}
Concretely, $r(z)$ is the meromorphic function of the spectral
parameter $z = z_1 - z_2$ obtained by extracting the OPE
against the Arnold form $\eta_{12} = d\log(z_1 - z_2)$ on
the two-point ordered configuration space.
\end{definition}

\begin{remark}[Pole absorption]
\label{rem:pole-absorption}
The $r$-matrix has poles one order lower than the OPE, because
the logarithmic form $d\log(z)$ absorbs one pole: the OPE of two
fields with a pole of order~$p$ produces an $r$-matrix with a pole
of order $p-1$. This is the $d\log$ absorption rule: Heisenberg
OPE $\sim 1/(z-w)^2$ gives $r \sim 1/z$; Virasoro OPE
$\sim 1/(z-w)^4$ gives $r \sim 1/z^3$.
\end{remark}

\begin{example}[Heisenberg ordered bar]
\label{ex:heisenberg-bar}
Let $\cH_k$ be the Heisenberg chiral algebra at level~$k$,
generated by a single field $J$ of conformal weight $h = 1$. The
desuspended generator has degree $|s^{-1}J| = 0$. At degree~$2$ the
ordered bar component is
\[
  (s^{-1}\bar\cH_k)^{\otimes 2}
  \;\cong\; k \cdot [J|J],
\]
a one-dimensional space. The bar differential extracts the
collision residue of $J(z_1)J(z_2) \sim k/(z_1-z_2)^2$ against
$\eta = d\log(z_1-z_2)$, producing a single simple pole
$k/(z_1-z_2)$. The $r$-matrix is
\begin{equation}\label{eq:heis-r-matrix}
  r^{\mathrm{Heis}}(z) \;=\; k/z.
\end{equation}
At $k = 0$ the $r$-matrix vanishes: the algebra becomes
trivially commutative. At degree~$3$, the ordered differential squares
to zero by Stokes on the consecutive Stasheff faces of
$\overline{\FM}_3^{\mathrm{ord}}(\CC)$. The full Arnold relation
belongs to the holomorphic flatness and ordered-to-symmetric descent
calculation, not to the list of ordered boundary generators.
\end{example}

\begin{remark}[Heisenberg $R$-matrix convention: curved vs flat]\label{rem:heis-convention-eone-primacy}
The Heisenberg algebra $\cH_k$ admits a flat presentation
$r(z)=k/z$ and a curved presentation $R(z)=\exp(\hbar k/z)$. They
are related by $R(z)=\exp(\hbar r(z))$; equivalently, $r(z)$ is the
classical logarithm of the quantum monodromy.
\end{remark}

\begin{example}[Affine Kac--Moody ordered bar]
\label{ex:km-bar}
Let $\widehat{\fg}_k$ be the affine Kac--Moody algebra at level~$k$
for a simple Lie algebra $\fg$ of dimension $\dim(\fg)$ and dual
Coxeter number $h^\vee$. The OPE has both simple and double poles:
$J^a(z) J^b(w) \sim k\kappa^{ab}/(z-w)^2 + f^{ab}{}_c J^c(w)/(z-w)$,
where $\kappa^{ab}$ is the Killing form and $f^{ab}{}_c$ the
structure constants. The $r$-matrix in the trace-form convention is
\begin{equation}\label{eq:km-r-matrix}
  r^{\KM}(z) \;=\; k\,\Omega_\fg/z,
\end{equation}
where $\Omega_\fg = \sum_a t^a \otimes t_a \in \fg \otimes \fg$
is the inverse Killing form Casimir. The level prefix~$k$ is
essential: at $k = 0$ the $r$-matrix vanishes (the algebra becomes
abelian in the sense that the double pole, which governs the
$r$-matrix, is zero). At the critical level $k = -h^\vee$, the
Sugawara construction fails; the behaviour of the $r$-matrix at
this locus is controlled by the Feigin--Frenkel centre.

The KZ normalization
$r_{\mathrm{KZ}}(z)=\Omega_{\mathrm{KZ}}/((k+h^\vee)z)$ is a separate
spectral convention. It is not the same rational function of the
level as the trace-form collision residue \(k\Omega_{\mathrm{tr}}/z\);
the Drinfeld--Kohno comparison relates their monodromy categories
after the standard KZ normalization has been chosen.
\end{example}

\begin{example}[Virasoro ordered bar]
\label{ex:vir-bar}
Let $\Vir_c$ be the Virasoro algebra at central charge~$c$. The
OPE $T(z)T(w) \sim (c/2)/(z-w)^4 + 2T(w)/(z-w)^2 + \partial T(w)/(z-w)$
has a quartic pole. The $d\log$ absorption gives a cubic-plus-simple
$r$-matrix:
\begin{equation}\label{eq:vir-r-matrix}
  r^{\Vir}(z) \;=\; (c/2)/z^3 + 2T/z.
\end{equation}
This is cubic plus simple, not quartic.
\end{example}

% ================================================================
% SECTION 3
% ================================================================
\section{The $\Eone$ convolution algebra and its MC element}
\label{sec:e1-convolution}

\subsection{The ribbon modular operad}

The combinatorial unit of the commutative modular operad
$\cM_\Com$ is a stable graph with symmetric vertex labels: ordering
data is erased before composition starts. One modular operad retains
the ordering without breaking the genus decomposition: the ribbon
modular operad $\cM_\Ass$.

\begin{definition}[Ribbon modular operad {\cite{GK98,CG17}}]
\label{def:ribbon-modular-operad}
The \emph{ribbon modular operad} $\cM_\Ass$ has components
\begin{equation}\label{eq:ribbon-mod-operad}
  \cM_\Ass(g,n)
  \;:=\;
  C_*\!\bigl(\overline{\cM}_{g,n}^{\,\mathrm{rib}}\bigr),
  \qquad 2g-2+n > 0,
\end{equation}
where $\overline{\cM}_{g,n}^{\,\mathrm{rib}}$ classifies stable
genus-$g$ Riemann surfaces with $n$ parametrised boundary circles.
A ribbon structure on a stable graph is a cyclic ordering of
half-edges at each vertex. In genus~$0$:
$\cM_\Ass(0,n) = C_*(K_n)$, where $K_n$ is the Stasheff
associahedron of dimension $n-2$, and grafting is face inclusion.
The composition maps are self-sewing
$\xi_{ij}\colon \cM_\Ass(g,n) \to \cM_\Ass(g+1,n-2)$ and
grafting
$\circ_i\colon \cM_\Ass(g_1,n_1) \otimes \cM_\Ass(g_2,n_2)
\to \cM_\Ass(g_1+g_2,n_1+n_2-2)$.
\end{definition}

\begin{definition}[Feynman transform $\FMAss$]
\label{def:fass}
The \emph{associative Feynman transform} $\FMAss$ is the modular
cooperad with $(g,n)$-component
\begin{equation}\label{eq:fass-def}
  (\FMAss)(g,n)
  \;:=\;
  \bigoplus_{\Gamma \in \mathsf{Gr}^{\mathrm{rib}}_{g,n}}
  \frac{1}{|\Aut(\Gamma)|}\,
  \det\bigl(E(\Gamma)\bigr)
  \;\otimes\;
  \bigotimes_{v \in V(\Gamma)}
  \cM_\Ass\bigl(g(v),n(v)\bigr)^\vee,
\end{equation}
summed over connected ribbon graphs of genus~$g$ with $n$ legs.
The differential
$D_{\FMAss} = \sum_{e \in E(\Gamma)} \pm\, \xi_e^*$
is dual to edge contraction. This differential satisfies
$D_{\FMAss}^2 = 0$ by the standard codimension-$2$ cancellation
argument~\cite{GK98}.
\end{definition}

\begin{theorem}[Coinvariant projection {\cite[Thm.~4.9]{GK98}}]
\label{thm:fass-to-fcom}
The composite of the ribbon-forgetting and external
$\Sigma_n$-coinvariant projections
\[
  (\FMAss)(g,n)
  \;\xrightarrow{\;\pi_{\mathrm{rib}}\;}
  (\FMAss)(g,n)/{\sim_{\mathrm{rib}}}
  \;\xrightarrow{\;\pi_{\Sigma_n}\;}
  (F\!\Com)(g,n)
\]
is a quasi-isomorphism of modular cooperads.
\end{theorem}

\subsection{The $\Eone$ convolution dg~Lie algebra}

\begin{definition}[The $\Eone$ modular convolution algebra]
\label{def:e1-convolution}
Let $\cA$ be a cyclic $\Eone$-chiral algebra. The
\emph{$\Eone$ modular convolution dg~Lie algebra} is
\begin{equation}\label{eq:e1-conv-def}
  \gEone
  \;:=\;
  \prod_{\substack{g,n \\ 2g-2+n > 0}}
  \Hom\!\bigl(
    \cM_\Ass(g,n),\,
    \End_\cA(n)
  \bigr).
\end{equation}
The Hom carries \emph{no} $\Sigma_n$-equivariance. This is the
structural distinction from the modular convolution algebra
\begin{equation}\label{eq:mod-conv}
  \gmod
  \;=\;
  \prod_{\substack{g,n \\ 2g-2+n > 0}}
  \Hom_{\Sigma_n}\!\bigl(
    \cM_\Com(g,n),\,
    \End_\cA(n)
  \bigr),
\end{equation}
whose Hom-source uses the symmetric cooperad $\Sym^c$ obtained
from $T^c$ by external $\Sigma_n$-coinvariants.
The dg~Lie structure on $\gEone$ is inherited from $\FMAss$:
the differential~$D$ comes from $D_{\FMAss}$ (edge contraction),
and the bracket $[-,-]$ comes from ribbon-graph composition.
\end{definition}

\begin{theorem}[$\Eone$ Maurer--Cartan element]
\label{thm:e1-mc}
The element
\begin{equation}\label{eq:theta-e1}
  \Theta^{\Eone}_\cA
  \;:=\;
  D^{\Eone}_\cA - d_0
  \;\in\; \gEone,
\end{equation}
where $D^{\Eone}_\cA$ is the ordered genus-completed bar
differential and $d_0$ the background differential, satisfies
the Maurer--Cartan equation
\begin{equation}\label{eq:e1-mc-equation}
  D\,\Theta^{\Eone}_\cA
  + \tfrac{1}{2}[\Theta^{\Eone}_\cA,
    \Theta^{\Eone}_\cA] = 0
\end{equation}
in the ordered dg~Lie algebra $\gEone$.
\end{theorem}

\begin{proof}
The identity $(D^{\Eone}_\cA)^2 = 0$ is equivalent to the MC
equation~\eqref{eq:e1-mc-equation} upon writing
$D^{\Eone}_\cA = d_0 + \Theta^{\Eone}_\cA$. The square-zero
property is the ordered bar-intrinsic construction:
$D^{\Eone}_\cA$ is built from the associative Feynman transform,
and $D_{\FMAss}^2 = 0$ forces the induced coderivation to square
to zero.
\end{proof}

\begin{proposition}[$\Eone$ shadow at each degree]
\label{prop:e1-shadow}
The degree-$r$ component of $\Theta^{\Eone}_\cA$ at genus~$0$
is an element
$r_r \in \End(V^{\otimes r}) \otimes H^{r-2}(K_r)
\otimes \cO(*\Delta)$,
whose $\Sigma_r$-coinvariant is the $\Einf$ shadow at degree~$r$:
\begin{center}
\renewcommand{\arraystretch}{1.15}
\begin{tabular}{lcc}
\toprule
\emph{Degree} & \emph{$\Eone$ shadow} & \emph{$\Einf$ shadow} \\
\midrule
$r = 2$ & $r(z)$ (classical $r$-matrix)
  & $\kappa(\cA)$ \\
$r = 3$ & $r_3(z_1,z_2)$ (associator shadow)
  & $\mathfrak{C}(\cA)$ \\
$r = 4$ & $r_4(z_1,z_2,z_3)$ (quartic shadow)
  & $\mathfrak{Q}(\cA)$ \\
\bottomrule
\end{tabular}
\end{center}
The MC equation at degree~$r$ is the $\partial K_{r+1}$-identity:
classical Yang--Baxter at $r = 2$ (two faces of $K_3$), pentagon
at $r = 3$ (five faces of $K_4$), and the quartic $R$-matrix
identity at $r = 4$.
\end{proposition}

\begin{proof}
At genus~$0$, $\cM_\Ass(0,r) = C_*(K_r)$. The MC equation
projected to degree~$r$ is the equation on the boundary of $K_{r+1}$,
whose faces are precisely the binary trees contributing to the
$r$-point collision structure. At $r = 2$, $K_3$ is an interval with
two vertices (the two binary trees on three leaves), and the
boundary identity is the classical Yang--Baxter equation.
\end{proof}

% ================================================================
% SECTION 4
% ================================================================
\section{The averaging map: surjectivity, MC projection,
  non-splitting}
\label{sec:averaging}

\subsection{Definition of $\av$}

\begin{definition}[The \(R\)-twisted averaging map]
\label{def:averaging-map}
With the notation of Definition~\ref{def:e1-convolution}, and after a
strong-unitary \(R\)-twisted descent datum has been fixed, the
\emph{averaging map} is the dg~Lie morphism
\begin{equation}\label{eq:av-def}
  \av_R\colon \gEone \twoheadrightarrow \gmod,
  \qquad
  \av_R(\phi)(g,n)
  \;:=\;
  \frac{1}{n!}\sum_{\sigma \in \Sigma_n}
  \sigma \cdot
  \bigl(\phi(g,n) \circ \iota^{\mathrm{rib}}_{g,n}\bigr),
\end{equation}
where
$\iota^{\mathrm{rib}}_{g,n}\colon \cM_\Com(g,n)
\to \cM_\Ass(g,n)$
is any section of the ribbon-forgetting quotient
(Theorem~\ref{thm:fass-to-fcom}). The result is
$\Sigma_n$-invariant by construction and independent of the choice
of section up to boundary terms. This is the Reynolds operator for
the external $\Sigma_n$-action: it takes external coinvariants
after pullback along~$\iota^{\mathrm{rib}}$.
\end{definition}

\begin{theorem}[$\av_R$ is a surjective dg~Lie morphism]
\label{thm:av-surjective}
Under the fixed strong-unitary \(R\)-twisted descent datum, the
averaging map
$\av_R\colon \gEone \twoheadrightarrow \gmod$
is a surjective morphism of dg~Lie algebras. Without this descent
datum there is only the ordered dg Lie algebra; no morphism to
\(\gmod\) is asserted.
\end{theorem}

\begin{proof}
Surjectivity: every $\Sigma_n$-invariant homomorphism from
$\cM_\Com(g,n)$ lifts, via any section of the ribbon-forgetting
quotient, to a homomorphism from $\cM_\Ass(g,n)$.
Compatibility with the differential: the ribbon edge-contraction
differential descends to the symmetric edge-contraction differential
under \emph{$R$-twisted} external $\Sigma_n$-coinvariants
$(-)^{R\text{-}\Sigma_n\text{-coinv}}$ --- the twist is trivial on the
pole-free subclass ($R=\id$), reducing to naive coinvariants, and
carries the Heisenberg monodromy $R(z) = \exp(k\hbar/z)$ on curved
families. Under this twist the comparison between ribbon and
commutative Feynman transforms~\cite{GK98} gives the differential
compatibility. Compatibility with the Lie bracket:
symmetrisation commutes with the convolution bracket on cooperadic
Homs because operad composition in $\FMAss$ is
$\Sigma_n$-equivariant.
\end{proof}

\subsection{Degree-$2$ averaging: the $r$-matrix projects to
$\kappa$}

\begin{theorem}[Degree-$2$ averaging]
\label{thm:degree-2-av}
Let $r(z) \in \End(V \otimes V) \otimes \cO(*\Delta)$ be the
genus-$0$ degree-$2$ $\Eone$ shadow of $\cA$ (the classical
$r$-matrix in the pre-dualisation convention). Then:
\begin{equation}\label{eq:degree-2-av}
  \av_{n=2}\bigl(r(z)\bigr) \;=\;
  \begin{cases}
    \kappa(\cA), &
      \text{for abelian Heisenberg families,} \\[4pt]
    \kappa_{\mathrm{dp}}(\widehat{\fg}_k)
    \;:=\; \dfrac{k\,\dim(\fg)}{2h^\vee}, &
      \text{for non-abelian affine KM}\\
      & \text{in the trace-form convention.}
  \end{cases}
\end{equation}
In the affine Kac--Moody case the full modular characteristic is
\begin{equation}\label{eq:km-kappa-full}
  \kappa(\widehat{\fg}_k)
  \;=\;
  \kappa_{\mathrm{dp}}(\widehat{\fg}_k) + \frac{\dim(\fg)}{2}
  \;=\;
  \frac{\dim(\fg)(k+h^\vee)}{2h^\vee},
\end{equation}
where the extra $\dim(\fg)/2$ is the Sugawara simple-pole
self-contraction through the adjoint Casimir eigenvalue $2h^\vee$.
\end{theorem}

\begin{proof}
The degree-$2$ MC component is, by definition, an element
$r(z) \in \End(V^{\otimes 2}) \otimes \cO(*\Delta)$ in the
pre-dualisation convention. Averaging over $\Sigma_2$ projects onto
the trivial isotypic component. For abelian Heisenberg families
this scalar coincides with the full modular characteristic. For
non-abelian affine Kac--Moody in the trace-form convention
$r(z) = k\,\Omega_\fg/z$, the averaging sees only the double-pole
channel and yields
$\kappa_{\mathrm{dp}} = k\,\dim(\fg)/(2h^\vee)$; the simple-pole
self-contraction contributes the additional Sugawara term
$\dim(\fg)/2$, so
$\kappa = \kappa_{\mathrm{dp}} + \dim(\fg)/2$.
The same formula is forced by pole-order analysis and by the limiting
case \(k=0\), where the trace-form \(r\)-matrix vanishes but the
Sugawara simple-pole contribution remains \( \dim(\fg)/2 \).
\end{proof}

\begin{remark}[Concrete examples]
\label{rem:concrete-r-kappa}
The principal families:
\begin{itemize}
\item Heisenberg $\cH_k$:
  $r(z) = k/z$,
  $\kappa(\cH_k) = k$.

\item Affine Kac--Moody $\widehat{\fg}_k$:
  $r(z) = k\,\Omega_\fg/z$
  (the level~$k$ survives $d\log$ absorption; at $k=0$ the
  $r$-matrix vanishes identically),
  $\av(r(z)) = \kappa_{\mathrm{dp}}
  = k\,\dim(\fg)/(2h^\vee)$,
  and
  $\kappa(\widehat{\fg}_k)
  = \dim(\fg)(k+h^\vee)/(2h^\vee)$.

\item Virasoro $\Vir_c$:
  $r(z) = (c/2)/z^3 + 2T/z$
  (cubic plus simple, not quartic),
  and $\kappa(\Vir_c) = c/2$.
\end{itemize}
\end{remark}

\subsection{Degree-$3$ averaging: the associator projects to the
cubic shadow}

\begin{theorem}[Degree-$3$ averaging]
\label{thm:degree-3-av}
Let $\Phi_{\mathrm{KZ}}(\cA) \in \End(V^{\otimes 3})$ be the
genus-$0$ degree-$3$ Maurer--Cartan component, which for affine
Kac--Moody $\widehat{\fg}_k$ is the Drinfeld--Kontsevich KZ
associator. Then
\begin{equation}\label{eq:degree-3-av}
\av_{R,n=3}\bigl(\Phi_{\mathrm{KZ}}(\cA)\bigr)
  \;=\; \mathfrak{C}(\cA),
\end{equation}
where $\mathfrak{C}(\cA)$ is the cubic shadow of Volume~I,
equivalently the $\Sigma_3$-coinvariant of the associator. The
antisymmetric component $[\Omega_{12}, \Omega_{23}]$ lies entirely
in $\ker(\av_R)$: it is antisymmetric under the transposition
$(1 \leftrightarrow 3)$, so its $\Sigma_3$-average vanishes. Only
the symmetric products
$\Omega_{ij}\Omega_{k\ell} + \Omega_{k\ell}\Omega_{ij}$ contribute
to~$\mathfrak{C}(\cA)$.
\end{theorem}

\begin{proof}
The $\Sigma_3$-isotypic decomposition of $\Phi_{\mathrm{KZ}}$ in
$\End(V^{\otimes 3})$ splits into the trivial, sign, and standard
representations. The Reynolds operator
$\av_R = (1/6)\sum_{\sigma \in \Sigma_3} R_\sigma\sigma$ projects onto the
trivial isotypic component. A direct check on
$[\Omega_{12}, \Omega_{23}]$ shows it is antisymmetric under
$(1 \leftrightarrow 3)$ and therefore has trivial isotypic component
zero. The identification of the coinvariant with the cubic shadow
$\mathfrak{C}(\cA)$ is~\cite[Theorem~D(iii)]{Lor26}.
\end{proof}

\subsection{MC projection}

\begin{theorem}[MC projection under $\av_R$]
\label{thm:av-mc-projection}
Under \(R\)-twisted averaging:
\begin{equation}\label{eq:av-mc}
  \av_R\bigl(\Theta^{\Eone}_\cA\bigr)
  \;=\; \Theta_\cA.
\end{equation}
In particular, the shadow obstruction tower
$(\kappa, \mathfrak{C}, \mathfrak{Q}, \ldots)$ is the
degree-by-degree $\Sigma_n$-coinvariant image of the $R$-matrix
tower $(r(z), r_3, r_4, \ldots)$.
\end{theorem}

\begin{proof}
$\av_R(D^{\Eone}_\cA) = D_\cA$ because the ordered genus-completed
bar differential on $\Barord(\cA)$ descends to $\BarSig(\cA)$
under $R$-twisted $\Sigma_n$-coinvariants, and
$\av_R(d_0) = d_0$. On the pole-free subclass the twist is trivial
and descent reduces to naive coinvariants; for Heisenberg, affine
Kac--Moody, and higher $E_\infty$-chiral families with poles, the
twist carries the monodromy $R(z)$ along elementary transpositions,
and the descent remains well-defined by quantum Yang--Baxter.
Therefore
\[
  \av_R(\Theta^{\Eone}_\cA)
  = \av_R(D^{\Eone}_\cA - d_0)
  = D_\cA - d_0
  = \Theta_\cA. \qedhere
\]
\end{proof}

\subsection{Non-splitting}

\begin{proposition}[Non-splitting of $\av_R$]
\label{prop:nonsplitting}
The short exact sequence of dg~Lie algebras
\begin{equation}\label{eq:nonsplitting-ses}
  0 \;\longrightarrow\; \ker(\av_R) \;\longrightarrow\;
  \gEone \;\xrightarrow{\;\av_R\;}\; \gmod
  \;\longrightarrow\; 0
\end{equation}
does not split as dg~Lie algebras for non-abelian families with a
nonzero degree-$3$ associator. Four independent obstructions prevent
a section:
\begin{enumerate}[label=\textup{(\roman*)}]
\item \emph{Fixed-degree triviality.}
  At each fixed degree~$n$, the inclusion section
  $s\colon \End^{\Sigma_n}(V^{\otimes n})
  \hookrightarrow \End(V^{\otimes n})$
  is a Lie subalgebra embedding, so the extension $2$-cocycle
  vanishes. The extension is Lie-trivial at each degree separately.

\item \emph{Cross-degree obstruction.}
  The bar differential
  $D\colon \gEone(n) \to \gEone(n{-}1)$
  contracts an internal edge, reducing degree by~$1$. This
  contraction is $\Sigma_{n-2}$-equivariant but
  \emph{not} $\Sigma_n$-equivariant: it breaks the symmetry from
  $\Sigma_n$ to $\Sigma_{n-2}$, leaking from
  $\ker(\av_{R,n})$ into $\operatorname{im}(\av_{R,n-1})$. No
  Lie-algebra section can simultaneously respect the fixed-degree
  splitting and intertwine with~$D$.

\item \emph{Associator representative.}
  At degree~$3$, the linearised Drinfeld associator
  $[\Omega_{12}, \Omega_{23}] \in \End(V^{\otimes 3})$ lies
  entirely in $\ker(\av_R)$: it is antisymmetric under
  $(1 \leftrightarrow 3)$, so its $\Sigma_3$-average vanishes.
  The infinitesimal braid relation
  $[\Omega_{12}, \Omega_{13} + \Omega_{23}] = 0$
  ensures that the commutator is intrinsic, not an artefact
  of the representation.

\item \emph{$\GRT_1$-torsor structure.}
  The space of dg~Lie sections of $\av_R$ satisfying the MC equation
  up to gauge is a torsor for the pro-unipotent
  Grothendieck--Teichm\"uller group $\GRT_1$ of
  Drinfeld~\cite{Drinfeld91}. This is the algebraic shadow of the
  Etingof--Kazhdan~\cite{EK96} non-uniqueness of quantisation: a
  splitting would canonically reconstruct the quantum group from the
  modular shadow $\kappa$ alone.
\end{enumerate}
For Heisenberg ($\dim V = 1$): $\ker(\av_R) = 0$ at every degree,
the extension splits trivially, and no $\GRT_1$ ambiguity arises.
\end{proposition}

\begin{proof}
Part~(i): the commutant of $\Sigma_n$ in $\End(V^{\otimes n})$ is
a Lie subalgebra, so $c(A,B) = (1 - \av_R)([A,B]) = 0$.

Part~(ii): the partial-trace map
$D_{ij}\colon \End(V^{\otimes n}) \to \End(V^{\otimes(n-1)})$
satisfies
$D_{ij} \circ \sigma = \sigma|_{\{1,\ldots,n\}\setminus\{j\}}
\circ D_{ij}$
only for $\sigma$ fixing both $i$ and $j$, i.e.\
$\Sigma_{n-2} \subset \Sigma_n$. A $\Sigma_n$-invariant
endomorphism in $\ker(\av_{R,n})$ can map under $D$ to an
endomorphism with nonzero $\av_{R,n-1}$ component.

Part~(iii): the $\Sigma_3$-isotypic decomposition of
$[\Omega_{12}, \Omega_{23}]$ consists of sign and standard
components only, with zero trivial component (verified in
$\End((\CC^2)^{\otimes 3})$). The infinitesimal braid relation
$[\Omega_{12}, \Omega_{13} + \Omega_{23}] = 0$ ensures the result
is intrinsic.

Part~(iv) follows from Drinfeld's theorem~\cite{Drinfeld91} on the
$\GRT_1$-torsor structure of associators, together with the
Etingof--Kazhdan quantisation theorem~\cite{EK96}.
\end{proof}

\begin{proposition}[Non-splitting at genus~$1$]
\label{prop:nonsplitting-genus1}
The short exact sequence~\eqref{eq:nonsplitting-ses}
restricted to genus~$1$ does not split. The three genus-$0$
obstructions persist, and a fourth obstruction arises from the
quasi-modular anomaly of the Eisenstein series~$E_2$. The
genus-$1$ propagator on the elliptic curve~$E_\tau$ is
$\eta_{12}^{(1)} = (\vartheta_1'/\vartheta_1)(z_1 - z_2 \mid
\tau)\,(dz_1 - dz_2)$. The self-energy regularisation produces the
Eisenstein series $E_2(\tau)$. The averaging map performs the
non-holomorphic completion
\[
  E_2(\tau)
  \;\longmapsto\;
  \widehat{E}_2(\tau)
  \;:=\;
  E_2(\tau) - \frac{3}{\pi\,\operatorname{Im}(\tau)}.
\]
A section $s\colon \gmod(1,1) \to \gEone(1,1)$ would lift
$\widehat{E}_2$ back to a holomorphic object on the $\Eone$ side;
the non-holomorphic correction $-3/(\pi\,\operatorname{Im}(\tau))$
has no preimage in the holomorphic $\Eone$ convolution Lie algebra,
so no such section exists.
\end{proposition}

\begin{proof}
The genus-$0$ obstructions persist via the embedding of the boundary
stratum $\overline{\cM}_{0,n+2} \subset \partial\overline{\cM}_{1,n}$.
The quasi-modular obstruction is independent: it is visible at
degree~$1$ (where $\Sigma_1$ is trivial and all genus-$0$
obstructions vanish), and it concerns the genus-$1$ fibre
(dependence on~$\tau$), not the degree (dependence on~$n$). The
two obstruction mechanisms are linearly independent in the relevant
$\operatorname{Ext}$ group.
\end{proof}

\begin{remark}[The kernel dimension formula]
\label{rem:ker-av-formula}
The kernel of $\av$ at degree~$n$ on the bare tensor power
$V^{\otimes n}$ has a closed-form dimension depending only on
$d = \dim V$:
\begin{equation}\label{eq:ker-av-dim}
  \dim\bigl(\ker(\av_n)\bigr)
  \;=\;
  d^n - \binom{n+d-1}{d-1}.
\end{equation}
The first term $d^n = \dim(V^{\otimes n})$ and the second
$\binom{n+d-1}{d-1} = \dim(\Sym^n V)$ is the dimension of the
$\Sigma_n$-coinvariant. This formula depends only on $\dim V$, not on
the Lie-algebraic structure of the representation, by the Schur--Weyl
analysis of~\cite[Proposition~6.1]{Lor26}. At $d = 1$
(Heisenberg): $\ker(\av_n) = 0$ for all~$n$. At $d = 2$
($\mathfrak{sl}_2$): $\ker(\av_2) = 1$,
$\ker(\av_3) = 4$, $\ker(\av_4) = 11$.
\end{remark}

% ================================================================
% SECTION 5
% ================================================================
\section{The $\Eone$ primacy theorem}
\label{sec:e1-primacy}

\subsection{Precise statement}

\begin{theorem}[$\Eone$ primacy]
\label{thm:e1-primacy-main}
Let $\cA$ be a cyclic chiral algebra on~$X$.
\begin{enumerate}[label=\textup{(\roman*)}]
\item \textup{(Surjectivity.)}
  The averaging map
  $\av_R\colon \gEone \twoheadrightarrow \gmod$
  is a surjective dg~Lie morphism after a strong-unitary
  \(R\)-twisted descent datum has been fixed.

\item \textup{(MC projection.)}
  $\av_R$ sends
  $\Theta^{\Eone}_\cA \mapsto \Theta_\cA$.
  In particular, the shadow obstruction tower
  $(\kappa, \mathfrak{C}, \mathfrak{Q}, \ldots)$
  is the degree-by-degree $\Sigma_n$-coinvariant image of the
  $R$-matrix tower $(r(z), r_3, r_4, \ldots)$.

\item \textup{(Bracket preservation.)}
  $\av_R$ preserves the dg~Lie bracket by
  $\Sigma_n$-equivariance of the operad composition maps.

\item \textup{(Non-splitting.)}
  For non-abelian families with a non-zero degree-$3$ associator
  component, the non-trivial non-coinvariant component represented
  infinitesimally by $[\Omega_{12},\Omega_{23}]$ lies in
  $\ker(\av_R)$. Its trivial isotypic projection is the cubic shadow
  $\mathfrak C(\cA)$; this scalar shadow does not reconstruct
  $\Phi_{\mathrm{KZ}}$. The short
  exact sequence
  $0 \to \ker(\av_R) \to \gEone
  \xrightarrow{\av_R} \gmod \to 0$
  does not split as dg~Lie algebras in these non-abelian cases. In
  the rank-one Heisenberg case $\ker(\av_n)=0$ for all $n$, so the
  degreewise extension splits trivially and no $\GRT_1$ ambiguity
  appears. At genus~$1$, a fourth
  obstruction to splitting arises from the quasi-modular anomaly
  of $E_2$ (Proposition~\textup{\ref{prop:nonsplitting-genus1}}).
\end{enumerate}
\end{theorem}

\begin{proof}
Parts~(i)--(iii) are established in
Theorem~\ref{thm:av-surjective} (surjectivity and bracket
preservation as a dg~Lie morphism) and
Theorem~\ref{thm:av-mc-projection} (MC projection
$\av_R(\Theta^{\Eone}_\cA) = \Theta_\cA$). The map $\av_R$ is the
Reynolds operator
$\phi \mapsto (1/n!)\sum_\sigma \sigma \cdot \phi$; surjectivity
follows because $\av$ restricts to the identity on
$\Sigma_n$-invariant elements, which span the target~$\gmod$.
The bracket is preserved because operad composition in $\FMAss$ is
$\Sigma_n$-equivariant: coinvariant projection converts
$T^c$-convolution to $\Sym^c$-convolution
(Theorem~\ref{thm:fass-to-fcom}).

For~(iv), the antisymmetric linearised associator
$[\Omega_{12},\Omega_{23}] \in \ker(\av_R)_{0,3}$ satisfies the
linearised genus-$0$, degree-$3$ MC equation in $\gEone$. The full
associator has a trivial-isotypic cubic shadow
$\mathfrak C(\cA)=\av_R(\Phi_{\mathrm{KZ}})$, but that scalar shadow
does not determine the ordered associator. Were $\av_R$ split, there
would exist a dg~Lie section
$s\colon \gmod \hookrightarrow \gEone$ with $\av_R \circ s = \id$.
Such a section would lift the commutative MC element $\Theta_\cA$
to an ordered element in the image of~$s$; but the pentagon
constraint on $\ker(\av_R)$ forces the degree-$3$ component of any
lift to involve $\Phi_{\mathrm{KZ}}$, which does not lie in the
image of any linear section
(Proposition~\ref{prop:nonsplitting}).
\end{proof}

\begin{principle}[The $\Eone$ primacy thesis]
\label{princ:e1-primacy}
Theorem~\textup{\ref{thm:e1-primacy-main}} admits a concise
restatement. The $\Eone$/ordered structure is the natural primitive
of the theory. The $\Einf$/symmetric/modular structure is its
$\Sigma_n$-coinvariant shadow under $\av_R$. The five main theorems
A--H are the invariants that survive averaging. Information flows
in only one direction: from the rich $\Eone$ side, which contains
the $R$-matrix, the KZ associator, the Drinfeld--Jimbo Yangian,
the full braided category of line operators, and the entire
spectral scattering data, to the coarser $\Einf$ side, which
contains only the modular characteristic $\kappa$, the shadow
obstruction tower, and the five structural theorems.
\end{principle}

\begin{remark}[Categorical level of the $\Etwo$ braided structure; ]
\label{rem:e2-lives-on-center}
The principle above states that the $\Eone$-side contains the
$R$-matrix and the braided category of line operators. To avoid the
most common conflation (MP1), we state explicitly \emph{which
object carries the $\Etwo$ structure}: the braided monoidal structure
is an $\Etwo$-algebra in $\mathrm{Cat}_\infty$ on
\emph{the Drinfeld center}
\[
  \cZ\!\bigl(\mathrm{Rep}^{\Eone}(\cA)\bigr),
\]
\emph{not} on $\cA$ itself. The algebra $\cA$ remains an $\Eone$-object;
the $\Etwo$ braiding arises one categorical level up, as the universal
half-braiding $\sigma_M(N)\colon M\otimes N\to N\otimes M$ constructed
from the $R$-matrix. Quantum groups $U_q(\fg)$ are the prototype: $U_q(\fg)$
is an $\Eone$-Hopf algebra, and
$\mathrm{Rep}(U_q(\fg))=\cZ(\mathrm{Rep}(\fg_0)_{\!\!\sqrt{\hbar}})$ is
$\Etwo$ in $\mathrm{Cat}$. The Deligne conjecture (classical and chiral)
upgrades this to an $\Etwo$-action on $\cZ_{\mathrm{ch}}^{\mathrm{der}}(\cA)
=C^\bullet_{\mathrm{ch}}(\cA,\cA)$; the cohomological version is
unconditional, the chain-level Tamarkin-type statement is conditional
on formality of $\mathrm{End}^{\mathrm{ch}}_\cA$.

In particular, Dunn additivity $\Eone\otimes\Eone=\Etwo$ applies to
$\cZ^{\mathrm{der}}_{\mathrm{ch}}(\cA)$ and $\mathrm{Mod}_\cA$, never
to $\cA$ itself. Compositions of the form ``$R$-matrix makes $\cA$ into
an $\Etwo$-algebra'' are category errors.
\end{remark}

\subsection{The Swiss-cheese / chiral Deligne--Tamarkin theorem}
\label{subsec:swiss-cheese-chiral-DT}

The dimensional uplift from a $2$-real-dimensional chiral algebra
on a curve $X$ to a $3$-real-dimensional holomorphic--topological
theory on $X \times \RR_{\geq 0}$ is not a property of the bar
direction alone. It is the universal property of the chiral
Swiss-cheese pair, and the bar coalgebra $\Barord(\cA)$ is the
chain-level computational model used to present that pair.

\begin{theorem}[Swiss-cheese / chiral Deligne--Tamarkin]
\label{thm:swiss-cheese-chiral-DT-standalone}
\textup{(}\emph{quadrant} Open, \emph{presentation} two-coloured
Swiss-cheese pair, \emph{level} $1 \to 3$ of the Beilinson tower,
\emph{hypothesis package}: a local $\Ainf$-chiral algebra $\cA$ on
a tangential log curve $(X, D, \tau)$ with completed
conformal-weight topology and finite chiral residue at every
double-point collision; a two-coloured chiral Swiss-cheese operad
$\SCchtop$ with chiral colour and topological collar
colour~\textup{\cite{Voronov99,KS00}}; an inner conformal vector
controlling $\bar Q$-translations on the topological collar.\textup{)}
The universal local chiral Swiss-cheese pair
\[
  U(\cA) \;=\; \bigl(\,
    \HH^\bullet_{\mathrm{ch}}(\cA, \cA),\; \cA,\;
    \{\mu^{\mathrm{univ}}_{p;q}\}\,\bigr)
  \;=\;
  \bigl(\,Z^{\mathrm{der}}_{\mathrm{ch}}(\cA),\; \cA,\;
  \{\mu^{\mathrm{univ}}_{p;q}\}\,\bigr)
\]
is terminal in the category of local $\SCchtop$-pairs over~$\cA$:
the closed colour $\HH^\bullet_{\mathrm{ch}}(\cA, \cA)
= Z^{\mathrm{der}}_{\mathrm{ch}}(\cA)$ acts on the open colour
$\cA$ through the chiral Swiss-cheese mixed operations
$\mu_{p;q}$ \textup{(}$p$ closed inputs, $q$ open inputs\textup{)};
the open colour does not act back on the closed colour. For any
local $\SCchtop$-pair $(B, \cA, \{\mu_{p;q}\})$ over $\cA$ with
continuous mixed operations there is a unique morphism
$\Phi \colon (B, \cA, \mu) \to U(\cA)$ over $\id_\cA$.
\end{theorem}

\begin{proof}[Citation]
This is Volume~I, Chapter on the chiral derived centre, the chiral
analogue of the Voronov Swiss-cheese formality
theorem~\cite{Voronov99} and the Kontsevich--Soibelman Deligne
conjecture~\cite{KS00}. Concretely, $\Phi(b)$ is the chiral
Hochschild cochain
\[
  \Phi(b)_n(a_1, \ldots, a_n;\, \lambda_1, \ldots, \lambda_{n-1})
  \;:=\;
  (-1)^{|b|}\,
  \mu_{1;n}(b;\, a_1, \ldots, a_n;\,
    \lambda_1, \ldots, \lambda_{n-1})
\]
on $n - 1$ open-sector spectral variables.
\end{proof}

\begin{remark}[The Swiss-cheese pair is the explanatory mechanism;
the bar is computational]
\label{rem:bar-is-computational-not-uplift}
The $2$-real to $3$-real dimensional uplift is the universal
normal-collar direction of the codimension-one boundary condition
forced by the chiral Deligne--Tamarkin terminal-object construction
of Theorem~\ref{thm:swiss-cheese-chiral-DT-standalone}. It is not
produced by iterating the bar functor on $\cA$. The bar coalgebra
$\Barord(\cA) = T^c(s^{-1}\bar\cA)$ is $\Eone$-coassociative on
the curve geometry: a single primitive twisting coalgebra over the
chiral associative operad $(\mathrm{ChirAss})^!$, with
deconcatenation coproduct and bar differential constituting one
graded coalgebra structure --- not two colours of a Swiss-cheese
algebra. The bar serves as a chain-level computational model that
resolves $\cA$ for the purpose of computing the closed-colour
Hochschild cochain complex via
\[
  Z^{\mathrm{der}}_{\mathrm{ch}}(\cA)
  \;\simeq\;
  \mathrm{Hom}_{\mathrm{ch}}\bigl(\Barord(\cA),\, \cA\bigr).
\]
The averaging map $\av_R$ acts within the open colour $\cA$,
projecting the ordered bar onto the symmetric bar; it does not
construct or witness the closed colour. The open and closed colours
are linked only through the Swiss-cheese mixed operations of
Theorem~\ref{thm:swiss-cheese-chiral-DT-standalone}.

In particular, the slogans \emph{``$\Barord(\cA)$ is the bulk''},
\emph{``the bar direction explains the $2$d $\to$ $3$d HT lift''},
and \emph{``the bar is the open--closed transition''} are
category-level errors: the bulk is the closed colour
$Z^{\mathrm{der}}_{\mathrm{ch}}(\cA)$, the dimensional uplift is
the universal collar of the Swiss-cheese pair, and the open--closed
transition is the chiral Deligne--Tamarkin morphism, not the bar
functor.
\end{remark}

\begin{remark}[Where the dimensional uplift comes from in
$\SCchtop$]
\label{rem:uplift-mechanism-source}
The Swiss-cheese chain-level construction adapts the topological
Voronov pair~\cite{Voronov99} to the chiral setting: the open
colour records $E_1$-chiral operations on $\cA$ along the curve
$X$ (one complex direction, two real directions); the closed
colour records $E_2$-topological brace operations on
$Z^{\mathrm{der}}_{\mathrm{ch}}(\cA)$ in the transverse half-plane
(two real directions). The chiral Deligne--Tamarkin terminal-object
construction supplies a universal morphism between these two
sectors that is itself a chiral object, with the inner conformal
vector of the algebra controlling $\bar Q$-translations on the
topological collar. The total $\SCchtop$ structure has
$2 + 1 = 3$ real dimensions: the holomorphic chiral pair on the
boundary curve and the topological collar normal to it. This is
the mechanism: not Costello--Gwilliam factorisation, which is
ambient-only, and not Voronov topological Swiss-cheese, which lacks
the holomorphic colour --- but their fusion via chiral
Deligne--Tamarkin. The bar of $\Eone$ is again $\Eone$ (cofree
conilpotent over $T^c$); the Swiss-cheese pair $(\Etwo, \Eone)$ at
the chiral-topological level is the source of the uplift.
\end{remark}

\subsection{The five ordered theorems}

Under the ordered Koszul and \(R\)-twisted descent hypotheses stated in
each item below, the five theorems of modular Koszul duality admit
ordered $\Eone$ counterparts. Their symmetric forms are obtained by
\(\av_R\), and no projection is asserted when the descent datum is absent.

\begin{theorem}[Theorem $\mathrm{A}^{\Eone}$: ordered bar-cobar]
\label{thm:ordered-A}
\textup{(}\emph{quadrant} Open, \emph{presentation} bar / twisting
coalgebra, \emph{level} $1 \leftrightarrow 2$ of the Beilinson
tower, \emph{hypothesis package}: strongly admissible augmented
cyclic $\Eone$-chiral algebra $\cA$ on a smooth curve, admissible
$R$-twisted descent datum.\textup{)}
The genus-$g$ ordered bar complex
$\Barord(\cA)(g,n)$ is an $\FMAss$-coalgebra with differential
satisfying $d^2 = 0$, where $\FMAss$ is the Feynman transform of
the ribbon modular operad. The bar and cobar functors form an
adjunction
\[
  \Omega^{\mathrm{ch},(g)} \dashv \overline{B}^{\mathrm{ord},(g)}\colon
  \mathrm{Alg}_{\FMAss} \rightleftarrows
  \mathrm{CoAlg}_{\FMAss},
\]
and, after \(R\)-twisted descent, the coinvariant projection recovers
Theorem~A of Volume~I.
\end{theorem}

\begin{theorem}[Theorem $\mathrm{B}^{\Eone}$: ordered inversion]
\label{thm:ordered-B}
\textup{(}\emph{quadrant} Open, \emph{presentation} bar / cobar
inversion, \emph{level} $1 \to 3$ of the Beilinson tower,
\emph{hypothesis package}: $\Eone$ Koszul locus
\textup{(}PBW-completeness on the ordered bar
complex\textup{)}, propagated genus by genus via the MC
perturbation lemma.\textup{)}
Assume $\cA$ lies on the $\Eone$ Koszul locus, meaning the
PBW-completeness hypothesis holds for the ordered bar complex. Then
at each genus~$g$, the counit
\[
  \Omega^{\mathrm{ch},(g)}\!\bigl(\Barord(\cA)(g,\bullet)\bigr)
  \;\xrightarrow{\;\sim\;}\; \cA
\]
is a quasi-isomorphism (inversion, not Koszul duality), and the
$\Eone$ linear Koszul dual $\cA^{!,\Eone}$ obtained by linear
duality on a finite-dimensional model of $\Barord(\cA)$ is an
$\Eone$-chiral algebra at each such genus.
\end{theorem}

\begin{theorem}[Theorem $\mathrm{C}^{\Eone}$: ordered
complementarity]
\label{thm:ordered-C}
\textup{(}\emph{quadrant} Open, \emph{presentation} ordered
complementarity, \emph{level} $5$ per stratum of the Beilinson
tower, \emph{hypothesis package}: ordered Koszul locus, finite
perfect model for the ordered bar pairing, ordered chiral
Hochschild centre $Z^{\mathrm{der}, \Eone}_{\mathrm{ch}}(\cA)$
identified with the centre used by the strict ordered
tower.\textup{)}
At each genus~$g$ and degree~$n$ with $2g-2+n>0$, the ordered
complementarity takes the form
\begin{equation}\label{eq:ordered-C}
  Q^{\Eone}_{g,n}(\cA) + Q^{\Eone}_{g,n}(\cA^{!,\Eone})
  \;=\;
  H^*\!\bigl(\overline{\cM}_{g,n}^{\mathrm{rib}},\,
  Z^{\Eone}(\cA)\bigr),
\end{equation}
where $Z^{\Eone}(\cA)$ is the ordered chiral Hochschild centre
$Z^{\mathrm{der}, \Eone}_{\mathrm{ch}}(\cA)$, defined by the
ordered chiral Hochschild cochain complex, not the subspace of
$R$-matrix-fixed vectors. The scalar identity
$\kappa + \kappa^! = 0$ for Kac--Moody and free-field families
(and $\kappa + \kappa^! = 13$ for Virasoro) is the
$\Sigma_n$-coinvariant image of this complementarity in the
verified scalar examples.
\end{theorem}

\begin{theorem}[Theorem $\mathrm{D}^{\Eone}$: ordered degree-$2$
package]
\label{thm:ordered-D}
\textup{(}\emph{quadrant} Open, \emph{presentation} degree-$2$
shadow tower, \emph{level} $4$ on $\overline\cM_{g,n}$ of the
Beilinson tower, \emph{hypothesis package}: genus-refined
degree-$2$ projection of $\Theta^{\Eone}_\cA$ exists genus by
genus; cyclic admissibility on each stratum.\textup{)}
The formal ordered degree-$2$ shadow series
\begin{equation}\label{eq:ordered-D}
  R^{\Eone,\mathrm{bin}}(z;\hbar)
  \;=\;
  \sum_{g \geq 0} \hbar^{2g}\, r_g(z)
\end{equation}
obtained from the genus-refined degree-$2$ projection of the $\Eone$
Maurer--Cartan element is the $\Eone$ degree-$2$ characteristic
package of~$\cA$. It is universal, additive, antisymmetric under
opposite-duality, and satisfies
$\av(r_0(z)) = \kappa(\cA)$ on abelian families, while for
non-abelian affine Kac--Moody
$\av(r_0(z)) = \kappa_{\mathrm{dp}}(\cA)$ and
$\kappa(\cA) = \kappa_{\mathrm{dp}}(\cA) + \dim(\fg)/2$.
The coinvariant recovery
$\av(R^{\Eone,\mathrm{bin}}) = \sum \hbar^{2g}\, F_g^{(2)}$
is the scalar genus-$2$ modular characteristic series of
Theorem~D.
\end{theorem}

\begin{theorem}[Theorem $\mathrm{H}^{\Eone}$: ordered
chiral Hochschild]
\label{thm:ordered-H}
\textup{(}\emph{quadrant} Open, \emph{presentation} ordered chiral
Hochschild concentration, \emph{level} $3$ of the Beilinson tower,
\emph{hypothesis package}: $\cA$ complete and ordered Koszul;
admissible $R$-twisted descent datum; complete bimodule $M$ in the
same completion as $\cA$; class M chain-level statement requires
weight-completed / pro-object / $J$-adic ambient.\textup{)}
For every genus~$g$ and complete $\cA$-bimodule~$M$, the genus-$g$
ordered chiral Hochschild homology is computed coalgebraically by
the genus-$g$ ordered chiral coHochschild complex of the ordered bar
coalgebra $C = \Barord(\cA)$, with twisted differential involving
the curvature term $\kappa(\cA) \cdot \lambda_g$ in the
uniform-weight case \textup{(UNIFORM-WEIGHT)}. At genus $g \geq 2$
with multi-weight field content the scalar formula receives a
cross-channel correction $\delta F_g^{\mathrm{cross}}$
\textup{(ALL-WEIGHT $+ \delta F_g^{\mathrm{cross}}$)}. The
$\Sigma_n$-coinvariant recovers the symmetric chiral Hochschild
complex of Theorem~H.
\end{theorem}

\begin{remark}[Theorem~D is the cleanest example]
\label{rem:thm-d-cleanest}
Among the five, Theorem~D makes the $\Eone$ primacy thesis most
concrete. The scalar modular characteristic $\kappa$ is a single
number per family, and the family formulas are distinct: for $\cW_N$
the harmonic number is $H_N = \sum_{j=1}^N 1/j$ and must not be
copied from a different family. Its degree-$2$ lift is a meromorphic
function of a single variable, the collision coordinate~$z$. The
statement $\av(r(z)) = \kappa$ says this exactly for Heisenberg,
while for non-abelian affine Kac--Moody the same averaging recovers
only $\kappa_{\mathrm{dp}}$ and the full $\kappa$ adds the Sugawara
shift $\dim(\fg)/2$. The scalar shadow is still the
$\Sigma_2$-average of a function. This is a strict projection;
the full function $r(z)$ carries information about the braiding, the
quantum group, and the line-operator algebra that no scalar invariant
can see.
\end{remark}

\begin{remark}[What averaging forgets]
\label{rem:what-av-forgets}
The kernel of $\av_R$ contains several objects of independent interest.
At degree~$2$, $\ker(\av_R)$ is the antisymmetric part of $r(z)$; for
bosonic generators of even desuspended degree this vanishes
identically, and $\kappa$ recovers all of $r(z)$. For generators of
odd desuspended degree, the kernel is nontrivial and captures the
parity-dependent Eulerian weight structure on the bar complex. At
degree~$3$, $\ker(\av_R)$ contains the linearised Drinfeld commutator
$[\Omega_{12}, \Omega_{23}]$ and, by the non-splitting proposition,
obstructs a canonical reconstruction of the full associator
from~$\mathfrak{C}$. The $\GRT_1$-torsor structure of the splittings
of $\av_R$ is the operadic shadow of Etingof--Kazhdan non-uniqueness
of quantisation.
\end{remark}

% ================================================================
% SECTION 6
% ================================================================
\section{The formality bridge: $\Einf$ input}
\label{sec:formality-bridge}

The ordered bar complex is defined for every $\Eone$-chiral algebra,
but the symmetric bar exists only when the $\cD_X$-module descends
to the unordered Ran space. For $\Einf$-chiral input the descent
holds, and the question is: does the coinvariant projection lose
information? The answer at the cohomological level is no; at the
chain level the loss is precisely the $R$-matrix and the associator.

\begin{remark}[Formality for $\Einf$-chiral input]
\label{rem:formality-einf}
For families whose nontrivial $\Sigma_n$-isotypic summands of the
ordered bar complex are acyclic, the $\Eone$ refinement adds no new
data at the level of the trivial-isotypic bar \emph{cohomology}: the
natural comparison map
\[
  H^*(\Barord(\cA))_{\mathrm{triv}} \;\xrightarrow{\;\sim\;}\;
  H^*(\BarSig(\cA))
\]
is an isomorphism. In characteristic zero the coinvariant functor is
exact, so taking coinvariants introduces no higher group-cohomology
obstruction. The ordered bar still carries chain-level data: the
$R$-matrix, the KZ connection, and the nontrivial component of the
Drinfeld associator determine the braided monoidal structure on the
category of $\cA$-modules.
\end{remark}

\subsection{The formality theorem}

\begin{theorem}[Formality bridge for $\Einf$-chiral algebras]
\label{thm:formality-bridge-standalone}
Let $\cA$ be an $\Einf$-chiral algebra on~$X$ and $\Phi$ a
Drinfeld associator. Assume the nontrivial $\Sigma_n$-isotypic
summands of the ordered bar complex are acyclic. The
Kontsevich--Tamarkin operad formality quasi-isomorphism induces a
quasi-isomorphism from the trivial-isotypic ordered bar to the
symmetric bar,
\begin{equation}\label{eq:formality-bridge-qi}
  \Barord(\cA)_{\mathrm{triv}}
  \;\xrightarrow[\Phi]{\;\sim\;}\;
  \BarSig(\cA)
\end{equation}
that intertwines the deconcatenation coproduct on $T^c(s^{-1}\bar\cA)$
with the coshuffle coproduct on $\Sym^c(s^{-1}\bar\cA)$ at the
cohomological level.

In particular, the coinvariant projection
$\pi_{\Sigma_n}\colon
(s^{-1}\bar\cA)^{\otimes n} \twoheadrightarrow
(s^{-1}\bar\cA)^{\odot n}$
induces an isomorphism on bar cohomology:
\begin{equation}\label{eq:bar-cohom-formality}
  H^*\bigl(\Barord(\cA)\bigr)
  \;\xrightarrow{\;\cong\;}\;
  H^*\bigl(\BarSig(\cA)\bigr).
\end{equation}
The transferred $\Ainf$-structure on the cohomological model
$H^*(\cA)$ may carry non-trivial higher operations
$m_k^{\mathrm{tr}}$ for $k \geq 3$; what the formality bridge
guarantees is that ordered and symmetric complexes compute the
same cohomology.
\end{theorem}

\begin{proof}
The argument proceeds in four steps.

\medskip
\noindent\textit{Step~1: Koszul resolution of the
$\Sigma_n$-coinvariant.}
Working in characteristic zero, the group algebra
$k[\Sigma_n]$ is semisimple (Maschke's theorem). The
idempotent
\begin{equation}\label{eq:reynolds-idempotent}
  e_n
  \;:=\;
  \frac{1}{n!}\sum_{\sigma \in \Sigma_n} \sigma
  \;\in\; k[\Sigma_n]
\end{equation}
is the Reynolds operator projecting $V^{\otimes n}$ onto
$\Sym^n(V) = (V^{\otimes n})_{\Sigma_n}$, and the complement
$(1 - e_n)(V^{\otimes n})$ is the kernel of $\av_n$.
Semisimplicity gives a direct sum decomposition
\begin{equation}\label{eq:maschke-decomp}
  V^{\otimes n}
  \;\cong\;
  \Sym^n(V) \;\oplus\; \ker(\av_n)
\end{equation}
as $k[\Sigma_n]$-modules. This decomposition is compatible
with the bar differential by the $\Sigma_n$-equivariance of
the OPE for $\Einf$-chiral algebras.

\medskip
\noindent\textit{Step~2: Group cohomology vanishing.}
In characteristic zero,
\begin{equation}\label{eq:group-cohom-vanishing}
  H^i(\Sigma_n,\, M) \;=\; 0
  \qquad\text{for all } i > 0
  \text{ and all } k[\Sigma_n]\text{-modules } M.
\end{equation}
This vanishing follows from semisimplicity: the category
$k[\Sigma_n]\text{-}\mathrm{mod}$ has global dimension zero,
so $\mathrm{Ext}^i_{k[\Sigma_n]}(k, M) = 0$ for $i > 0$.
The vanishing ensures that the coinvariant projection
$\pi_{\Sigma_n}\colon \Barord(\cA)
\twoheadrightarrow \BarSig(\cA)$ is exact: no higher
$\Sigma_n$-cohomology obstructs the passage from ordered to
symmetric.

\medskip
\noindent\textit{Step~3: The quasi-isomorphism.}
Consider the coinvariant short exact sequence at each
degree~$n$:
\begin{equation}\label{eq:coinv-ses}
  0 \;\to\; \ker(\av_n)
  \;\to\; (s^{-1}\bar\cA)^{\otimes n}
  \;\xrightarrow{\;\pi_{\Sigma_n}\;}
  (s^{-1}\bar\cA)^{\odot n} \;\to\; 0.
\end{equation}
The bar differential $d_{\mathrm{bar}}$ preserves the
decomposition~\eqref{eq:maschke-decomp} because the
$\Einf$-chiral OPE is $\Sigma_n$-equivariant: the collision
residue $a_i \star a_{i+1}$ is symmetric in $a_i, a_{i+1}$
up to the graded sign from the Arnold form. The
decomposition~\eqref{eq:maschke-decomp} at the chain level,
combined with the cohomology vanishing
of~\eqref{eq:group-cohom-vanishing}, gives
\[
  H^*\bigl(\ker(\av_n), d_{\mathrm{bar}}\bigr)
  \;\hookrightarrow\;
  H^*\bigl(\Sigma_n, (s^{-1}\bar\cA)^{\otimes n}\bigr)
  \;=\; 0
  \quad\text{for } * > 0.
\]
The long exact sequence in cohomology then forces
$\pi_{\Sigma_n}$ to be a quasi-isomorphism.

\medskip
\noindent\textit{Step~4: Explicit quasi-isomorphism map.}
The Kontsevich--Tamarkin operad formality
(Kontsevich~\cite{Kontsevich99},
Tamarkin, see also~\cite{LV12}) provides a
quasi-isomorphism
\[
  \varphi_\Phi \colon
  C_*(\overline{\FM}_n(\CC))
  \;\xrightarrow{\;\sim\;}\;
  H^*(\overline{\FM}_n(\CC))
\]
depending on the Drinfeld associator~$\Phi$. Tensoring with
$(s^{-1}\bar\cA)^{\otimes n}$ and taking $R$-twisted $\Sigma_n$-coinvariants
(equivalently $\BarSig_n(\cA) \simeq (\Barord_n(\cA))^{R\text{-}\Sigma_n\text{-coinv}}$;
the twist reduces to the naive one on the pole-free subclass;)
gives the quasi-isomorphism~\eqref{eq:formality-bridge-qi}.
The dependence on~$\Phi$ is a $\GRT_1$-ambiguity: different
choices of associator yield homotopic maps.
\end{proof}

\begin{remark}[What the bridge preserves and what it loses]
\label{rem:bridge-preserves-loses}
The formality bridge preserves the bar \emph{cohomology}: homological
invariants such as the dimension of the ordered bar cohomology groups,
the Koszul property (PBW-degree concentration), and the five theorem
statements. It loses chain-level data: the $R$-matrix $r(z)$ is a
chain-level datum on the ordered side that becomes trivial
($\Sigma_2$-average $\to \kappa$) on the symmetric side. The
Drinfeld associator $\Phi_{\mathrm{KZ}}$ is a chain-level datum at
degree~$3$. Its antisymmetric linearised commutator lies in
\(\ker(\av_R)\), while its trivial-isotypic projection is the cubic
shadow \(\mathfrak C(\cA)\). The braided monoidal structure on the
category of $\cA$-modules is encoded in the chain-level ordered bar
complex and is invisible from the symmetric bar cohomology.
\end{remark}

% ================================================================
% SECTION 7
% ================================================================
\section{The formality failure: $\Eone$ input}
\label{sec:formality-failure}

For genuinely $\Eone$-chiral input (Yangians, Etingof--Kazhdan
quantum vertex algebras), the formality bridge of
Section~\ref{sec:formality-bridge} breaks down. The obstruction
is clean: the $\cD_X$-module underlying the OPE does not descend
to the symmetric quotient $X^{(n)}$, because the $R$-matrix
$S(z) \neq \id$ breaks the $\Sigma_n$-equivariance of the bar
differential. The ordered bar is the only bar that exists, and
the averaging map has no target.

\subsection{The theorem}

\begin{theorem}[Formality failure for genuinely $\Eone$-chiral
algebras]
\label{thm:formality-failure-standalone}
Let $\cA$ be a genuinely $\Eone$-chiral algebra
(tier~(c) of the $\Eone$/$\Einf$ hierarchy: the vertex
$R$-matrix $S(z) \neq \id$ is not derived from any $\Einf$-chiral
OPE). The formality bridge of
Theorem~\textup{\ref{thm:formality-bridge-standalone}} fails in
the following precise sense.
\begin{enumerate}[label=\textup{(\roman*)}]
\item \textup{(Operad formality is unconditional.)}
  The Kontsevich--Tamarkin formality of
  $C_*(\overline{\FM}_n(\CC))$ as an operad depends only on
  the operad structure, not on the specific algebra. It holds
  for $\Eone$-chiral input.

\item \textup{(Equivariance fails: the $\cD$-module does not
  descend.)}
  The factorization $\cD$-module
  $\cF_n^{\mathrm{ord}}$ on $\Conf_n^{\mathrm{ord}}(X)$
  is \emph{not} $\Sigma_n$-equivariant. The obstruction is the
  $R$-matrix: the monodromy of $\cF_n^{\mathrm{ord}}$ around
  the diagonal $\Delta_{ij}$ is
  \begin{equation}\label{eq:monodromy-obstruction}
    \mathrm{Mon}_{\gamma_{ij}}(\cF_n^{\mathrm{ord}})
    \;=\;
    S_{ij}(z)
    \;\neq\; \id,
  \end{equation}
  where $S_{ij}(z)$ is the braiding by the $R$-matrix on the
  $i$-th and $j$-th tensor factors. Since
  $S_{ij} \neq \id$, the $\cD$-module does not descend from
  $\Conf_n^{\mathrm{ord}}(X)$ to the symmetric quotient
  $\Conf_n(X)/\Sigma_n$; equivalently, the symmetric bar
  $\BarSig(\cA)$ does not exist as a factorization coalgebra on
  the unordered Ran space $\Ran(X)$.

\item \textup{(Dimension count: the Arnold algebra.)}
  By Arnold's theorem~\cite{Arnold69}, the cohomology algebra
  $H^*(\Conf_n(\CC))$ is the quotient of the exterior algebra on
  generators $\omega_{ij} = [d\log(z_i - z_j)]$ of degree~$1$
  ($1 \leq i < j \leq n$) by the Arnold relation
  $\omega_{ij}\omega_{jk} + \omega_{jk}\omega_{ki}
  + \omega_{ki}\omega_{ij} = 0$.
  Its Poincar\'e polynomial is
  \begin{equation}\label{eq:arnold-poincare}
    P_n(t)
    \;=\;
    \prod_{j=0}^{n-1}(1 + jt).
  \end{equation}
  Setting $t = 1$ gives $P_n(1) = n!$ as total dimension.
  At $n = 2$: $P_2(1) = 2$; at $n = 3$: $P_3(1) = 6$;
  at $n = 4$: $P_4(1) = 24$. The ordered chiral chain complex
  at degree~$n$ has dimension
  \begin{equation}\label{eq:ordered-dim}
    \dim \cC_n^{\mathrm{ord}}(D^\times, \cA)
    \;=\;
    (\dim \bar\cA)^n \cdot n!,
  \end{equation}
  where the factor $n!$ from the Arnold algebra measures the
  holomorphic data carried by the ordered complex beyond the
  topological Hochschild complex (which has dimension
  $(\dim \cA)^n$, without the $n!$ factor).

\item \textup{(Braid group representation.)}
  The fundamental group $\pi_1(\Conf_n(\CC)) = \PBr_n$ is the
  pure braid group on $n$ strands. The $\cD$-module
  $\cF_n^{\mathrm{ord}}$ gives a representation
  \begin{equation}\label{eq:braid-rep}
    \rho_n \colon
    \PBr_n
    \;\longrightarrow\;
    \Aut\!\bigl(\cC_n^{\mathrm{ord}}(D^\times, \cA)\bigr)
  \end{equation}
  by monodromy. By the Drinfeld--Kohno theorem, the local
  monodromy around $\Delta_{ij}$ is $\exp(2\pi i \cdot r_{ij})$,
  where $r_{ij}$ is the classical $r$-matrix acting on
  the $i$-th and $j$-th factors. For $\Einf$ input,
  $r_{ij}$ is scalar and the braid group acts through
  $\Sigma_n$ (the abelianisation); for genuinely $\Eone$ input,
  the braid representation is faithful and the ordered chiral
  homology carries the full braided monoidal structure.
\end{enumerate}
\end{theorem}

\begin{proof}
Part~(i): Kontsevich--Tamarkin operad formality depends on the
structure maps of the Fulton--MacPherson operad, not on the
algebra to which it is applied. The formality map
$\varphi_\Phi$ of~\eqref{eq:formality-bridge-qi} exists
unconditionally.

Part~(ii): the factorization $\cD$-module $\cF_n^{\mathrm{ord}}$
on $\Conf_n^{\mathrm{ord}}(X)$ is constructed from the OPE of
$\cA$. The OPE $Y(a,z)b$ and $Y(b,-z)a$ are related by the
$R$-matrix: $Y(a,z)b = S(z) \cdot Y(b,-z)a$. If
$S(z) \neq \id$, the transposition
$\sigma_{ij}\colon (z_i, z_j) \mapsto (z_j, z_i)$ acts on
$\cF_n^{\mathrm{ord}}$ by $S_{ij}(z_i - z_j) \neq \id$, and the
$\cD$-module does not descend to $X^{(n)}$. The ordered complex
retains holomorphic data from the spectral-parameter dependence
of~$S(z)$ that has no counterpart in the cyclic bar complex.

Part~(iii): the Poincar\'e polynomial
$P_n(t) = \prod_{j=0}^{n-1}(1 + jt)$ follows from
Arnold's computation of $H^*(\Conf_n(\CC))$~\cite{Arnold69}:
the generators $\omega_{ij}$ are of degree~$1$, the Arnold
relation kills the excess, and the Hilbert function is
determined by the recursion $P_n(t) = (1 + (n-1)t)\cdot P_{n-1}(t)$
with $P_1(t) = 1$.

Part~(iv): the Drinfeld--Kohno theorem
(Drinfeld~\cite{Drinfeld91}, Kohno~\cite{Kohno85})
identifies the monodromy representation of the KZ connection
with the quantum group $R$-matrix action, giving the braid
representation~\eqref{eq:braid-rep}.
For $\Einf$ input, the $R$-matrix is the Koszul-signed identity
$S(z) = (-1)^{|a||b|}\id$, so the braid group acts through
signs. For genuinely $\Eone$ input, the representation is
faithful: the Yang--Baxter equation on $r(z)$ ensures
consistency of the braid relations, and non-triviality of $S(z)$
ensures non-triviality of the representation.
\end{proof}

\subsection{The Yangian as paradigm}

\begin{remark}[$n!$ versus $1$]
\label{rem:n-factorial-vs-1}
The dimension count exposes the gap: the $\Eone$ convolution space
$\Hom(\cM_\Ass(g,n), \End_\cA(n))$ has dimension
$n! \cdot \dim\bigl(\Hom_{\Sigma_n}(\cM_\Com(g,n),
\End_\cA(n))\bigr)$ at each degree, by the multiplicity of the
regular representation in the coinvariant decomposition. For
$\Einf$ input the $n!$ factor is redundant: the data is
$\Sigma_n$-equivariant and the coinvariant projection is a
quasi-isomorphism. For $\Eone$ input the $n!$ factor is essential:
the Maurer--Cartan element $\Theta^{\Eone}_\cA$ uses all $n!$
orderings, the $R$-matrix cannot be recovered from its
$\Sigma_2$-coinvariant, and the symmetric bar $\BarSig(\cA)$
does not exist as a factorization coalgebra on the
\emph{unordered} Ran space (the $\cD$-module does not descend to
$X^{(n)}$).
\end{remark}

The paradigmatic example is the Yangian $Y_\hbar(\fg)$. Its
$R$-matrix $r_{\mathrm{YB}}(z) = \hbar\,\Omega_\fg/z$ is a
solution of the classical Yang--Baxter equation, and the bar
complex $\Barord(Y_\hbar(\fg))$ is an $\FMAss$-coalgebra whose
ordered Koszul comparison is controlled by the affine
Drinfeld--Kohno local system under the flatness hypotheses of
Proposition~\ref{prop:yangian-ordered-koszul-sa}. The symmetric bar
$\BarSig(Y_\hbar(\fg))$ does not exist because the Yangian's OPE
is not symmetric: the braiding $S(z) = 1 + \hbar\,P/z$ satisfies
$S(z) \neq S(-z)$. This is the cleanest demonstration of $\Eone$
primacy: the ordered bar is the only bar that can be formed, and
the averaging map has no target.

\begin{definition}[Ordered Koszul chiral algebra]
\label{def:ordered-koszul-chiral-algebra-sa}
An $\Eone$-chiral algebra $\cA$ is \emph{ordered Koszul} if
\textup{(i)} the universal twisting morphism
$\tau^{\ord}_\cA \colon \Barord(\cA) \to \cA$ is a chiral
Koszul morphism in the ordered sense \textup{(}the bar--cobar
counit $\Omega^{\ord}(\Barord(\cA)) \to \cA$ is a
quasi-isomorphism of $\Eone$-chiral algebras\textup{)}; and
\textup{(ii)} the PBW spectral sequence of $\Barord(\cA)$
collapses at $E_2$ in the ordered grading.
\end{definition}

\begin{proposition}[Yangian ordered-Koszul criterion]
\label{prop:yangian-ordered-koszul-sa}
Under Etingof--Kazhdan PBW flatness, Koszulity of the pure-braid
Orlik--Solomon algebra, and the Drinfeld--Kohno monodromy comparison,
$Y_\hbar(\fg)^{\mathrm{ch}}$ is ordered Koszul in the sense of
Definition~\textup{\ref{def:ordered-koszul-chiral-algebra-sa}}.
\end{proposition}

\begin{proof}
Present $Y_\hbar(\fg)^{\mathrm{ch}}$ in Drinfeld's second form as
generated by $\{\xi^{\pm}_{i,r}, h_{i,r}\}$ modulo the
loop--algebra relations; by Etingof--Kazhdan flatness
\textup{\cite{EK96}} the resulting filtration is PBW-flat over
$\hbar$. Koszulity of the pure-braid Orlik--Solomon algebra
(Shelton--Yuzvinsky~\textup{\cite{SheltonYuzvinsky1997}}) identifies
the ordered bar cohomology with the quadratic dual under the
Drinfeld--KZ bijection; $E_2$-collapse is the
Etingof--Kazhdan--Drinfeld flatness statement. The EK monodromy
provides the counit quasi-isomorphism witness at generic
$\hbar$. The symmetric bar has no target, so the ordered
presentation is the unique one.
\end{proof}

\begin{lemma}[$R$-twisted $\Sigma_n$ descent (unitary $R$)]
\label{lem:R-twisted-descent-sa}
Let $\cA$ be an ordered Koszul $\Eone$-chiral algebra with
$R$-matrix $R(z) = 1 + \hbar\,r(z) + O(\hbar^2)$ satisfying the
classical Yang--Baxter equation \emph{and} unitarity
$R(z)\,R^{op}(-z) = \id$, and let
$\Barord(\cA)$ be its ordered bar complex. The twisted
$\Sigma_n$-action on $\Barord(\cA)_n$, defined on a transposition
$\sigma_{i,i+1}$ by
$\sigma_{i,i+1}\cdot(a_1\otimes\cdots\otimes a_n)
= R_{i,i+1}(z_i-z_{i+1})\cdot(a_1\otimes\cdots\otimes a_{i+1}\otimes a_i\otimes\cdots\otimes a_n)$,
extends to an action of the pure braid group $\mathrm{PBr}_n$ by
Yang--Baxter; unitarity implements the transposition relation
$\sigma_{i,i+1}^2 = \id$ at the monodromy level, descending the
action to $\Sigma_n$. The quotient
$\Barord(\cA)_n/\Sigma_n$ under this $R$-twist is the symmetric bar
$\BarSig(\cA)_n$ whenever the latter exists \textup{(}i.e.\ on
the $\Einf$-factorization locus\textup{)}. The lossy descent
map $\av\colon\Barord(\cA) \to \BarSig(\cA)$ is precisely the
$R$-twisted Reynolds projector associated with this action.
Unitarity is automatic after scalar normalization for rational
Yangians $Y_\hbar(\fg)$ and for trigonometric
$U_q(\widehat{\fg})$ at generic~$q$, vacuous for $R(z)=\tau$, and
convention-dependent for elliptic Belavin/Felder $R$-matrices
(cf.\ Vol~I, Lemma~\texttt{lem:R-twisted-descent} and its
scope remark).
\end{lemma}

This lemma is the $\Eone$ analogue of the classical
Dunn--Lurie $\Einf$-descent and records that the gap between
$\Barord$ and $\BarSig$ is controlled not by an abelian
$\Sigma_n$-action but by a braid representation.

\begin{example}[Yangian $Y_\hbar(\mathfrak{sl}_2)$:
the braid representation]
\label{ex:yangian-braid}
For $\fg = \mathfrak{sl}_2$ and $V = \CC^2$ the fundamental
representation, the $R$-matrix in the fundamental
representation is $R(u) = u\,I_4 + \hbar\,P \in \End(\CC^2
\otimes \CC^2)$, where $P(e_i \otimes e_j) = e_j \otimes e_i$.
The braid group $\PBr_2 = \ZZ$ acts on
$\cC_2^{\mathrm{ord}}(D^\times, Y_\hbar(\mathfrak{sl}_2))$
by the monodromy operator
\[
  \rho_2(\gamma_{12})
  \;=\;
  \exp(2\pi i \cdot r_{12})
  \;=\;
  \exp\!\bigl(2\pi i\,\hbar\,\Omega_{\mathfrak{sl}_2}/z\bigr),
\]
which at $z = 1$ and generic~$\hbar$ has four distinct eigenvalues
on $V \otimes V = V_0 \oplus V_1$ (the singlet and triplet of
$\mathfrak{sl}_2$). The symmetric bar would project this to the
$\Sigma_2$-coinvariant $\Sym^2(V) = V_1$, losing the singlet
channel; the ordered bar retains both isospin channels.
\end{example}

% ================================================================
% SECTION 8
% ================================================================
\section{Heisenberg and Kac--Moody: complete $\Eone$ computation}
\label{sec:heis-km}

This section gives the complete degree-$2$ ordered computation
for the two simplest families and records the passage through $\av$
to the scalar shadow. The computations are self-contained and
illustrate every step of the $\Eone$ primacy thesis.

\subsection{Heisenberg}

For $\cH_k$ with OPE $J(z)J(w) \sim k/(z-w)^2$:
\begin{itemize}
\item $r(z) = k/z$ (equation~\eqref{eq:heis-r-matrix}).
\item $\av(r(z)) = \kappa(\cH_k) = k$
  (the $\Sigma_2$-average of $k/z$ against $d\log(z)$ extracts
  the residue $k$, which is the full modular characteristic).
\item $\ker(\av_2) = 0$ (one-dimensional representation).
\item Bar cohomology has its reduced generator
  $H^{-1}(\Barord(\cH_k)) = \CC \cdot s^{-1}J$ in bar degree~$1$.
  Koszul dual:
  $\cH_k^! = (\cH_k^{\mathrm{i}})^\vee$ with
  $\kappa(\cH_k^!) = -k$. Complementarity:
  $\kappa(\cH_k) + \kappa(\cH_k^!) = k + (-k) = 0$
  (Kac--Moody and free-field families have $\kappa + \kappa^! = 0$;
  Virasoro has $\kappa + \kappa^! = 13$; the value is
  family-dependent).
\item The ordered Yangian is trivial: $Y_\hbar(\CC) \cong \CC[\hbar]$.
  The KZ connection is scalar (no braiding). The $\GRT_1$ ambiguity
  is absent.
\end{itemize}

The Heisenberg computation makes the passage through $\av$ explicit.
The degree-$2$ ordered MC component is
\begin{equation}\label{eq:heis-mc-deg2}
  \Theta^{\Eone,2}_{\cH_k}
  \;=\;
  k \cdot [s^{-1}J | s^{-1}J] \otimes \frac{dz}{z}
  \;\in\; (s^{-1}\bar\cH_k)^{\otimes 2}
  \otimes \Omega^1(\Conf_2^{\mathrm{ord}}(\CC)),
\end{equation}
a single generator tensored with the Arnold form $\eta = dz/z$ on
the two-point ordered configuration. The averaging map acts by
the residue extraction:
\begin{equation}\label{eq:heis-av-computation}
  \av(\Theta^{\Eone,2}_{\cH_k})
  \;=\;
  \frac{1}{2\pi i}\oint_{\gamma_0}
  \frac{k}{z}\,dz
  \;=\; k
  \;=\; \kappa(\cH_k).
\end{equation}
The kernel $\ker(\av_2) = 0$ because $\dim V = 1$, so the ordered
and symmetric bars carry identical degree-$2$ data. At degree~$3$
and higher, the Arnold relation is automatically satisfied (no
non-abelian Casimirs exist), and the Heisenberg shadow tower
terminates at $S_2 = \kappa = k$: this is class~$G$, with
$\Delta = 8\kappa \cdot S_4 = 0$ (since $S_4 = 0$).

The KZ connection at $n$ points is scalar:
\begin{equation}\label{eq:heis-kz}
  \nabla^{\cH_k}_{\mathrm{KZ}}
  \;=\; d - k\sum_{i < j} \frac{dz_{ij}}{z_i - z_j},
\end{equation}
and its monodromy is abelian: the representation
$\rho_n\colon \PBr_n \to \CC^\times$ factors through
$\PBr_n \twoheadrightarrow \ZZ^{\binom{n}{2}}$, and each
generator $\gamma_{ij}$ acts by the scalar $e^{2\pi i k}$.
No Drinfeld associator appears.

\subsection{Affine Kac--Moody}

For $\widehat{\fg}_k$ with $\fg$ simple, $\dim(\fg) = d$,
$h^\vee$ the dual Coxeter number:
\begin{itemize}
\item $r(z) = k\,\Omega_\fg/z$
  (equation~\eqref{eq:km-r-matrix}; trace-form convention;
  at $k=0$ the $r$-matrix vanishes identically).

\item $\av(r(z)) = \kappa_{\mathrm{dp}}
  = k\,d/(2h^\vee)$ (double-pole channel).

\item $\kappa(\widehat{\fg}_k) = d(k+h^\vee)/(2h^\vee)$
  (equation~\eqref{eq:km-kappa-full}; at $k = 0$:
  $\kappa = d/2$, not zero; at $k = -h^\vee$: $\kappa = 0$).

\item $\ker(\av_2) = (d^2 - d(d+1)/2)$-dimensional: the
  antisymmetric part of $\Omega_\fg$ in $\fg \otimes \fg$.
\end{itemize}

The degree-$2$ ordered MC component is
\begin{equation}\label{eq:km-mc-deg2}
  \Theta^{\Eone,2}_{\widehat\fg_k}
  \;=\;
  \sum_{a} k\,t^a \otimes t_a \;\cdot\;
  [s^{-1}J^a | s^{-1}J^b] \otimes \frac{dz}{z}
  \;\in\; (s^{-1}\bar V)^{\otimes 2}
  \otimes \Omega^1(\Conf_2^{\mathrm{ord}}(\CC)),
\end{equation}
where $\{t^a\}$ is a basis of $\fg$ with dual basis $\{t_a\}$
under the Killing form, and $V$ denotes the vacuum module of
$\widehat\fg_k$. The $\Sigma_2$-average extracts the trace over
$\fg$:
\begin{equation}\label{eq:km-av-computation}
  \av(\Theta^{\Eone,2}_{\widehat\fg_k})
  \;=\;
  \frac{1}{2\pi i}\oint_{\gamma_0}
  k\,\frac{\mathrm{tr}(\Omega_\fg)}{z}\,dz
  \;=\;
  k \cdot \frac{d}{2h^\vee}
  \;=\; \kappa_{\mathrm{dp}},
\end{equation}
where $\mathrm{tr}(\Omega_\fg)
= \sum_a \kappa^{ab}\kappa_{ab} = d/(2h^\vee)$ in the
normalisation where $\kappa(t^a, t_a) = 1$. The full
modular characteristic adds the Sugawara simple-pole
self-contraction:
\begin{equation}\label{eq:km-sugawara-shift}
  \kappa(\widehat\fg_k)
  \;=\;
  \kappa_{\mathrm{dp}} + \frac{d}{2}
  \;=\;
  \frac{d(k + h^\vee)}{2h^\vee}.
\end{equation}
The extra $d/2$ arises from the adjoint Casimir
eigenvalue $2h^\vee$: the simple-pole OPE
$f^{ab}{}_c J^c(w)/(z-w)$ contributes through the
self-contraction $\sum_{a,c} f^{ac}{}_d f^{ad}{}_c = 2h^\vee
\cdot \delta^{cd}$, which after normalisation yields $d/2$.

\subsubsection{The KZ connection and its monodromy}

The $n$-point KZ connection on the punctured disk is
\begin{equation}\label{eq:km-kz-connection}
  \nabla_{\mathrm{KZ}}
  \;=\; d \;-\; \sum_{1 \leq i < j \leq n}
  \frac{\Omega^{\mathrm{KZ}}_{ij}}{(k+h^\vee)(z_i - z_j)}\, dz_{ij},
\end{equation}
where \(\Omega^{\mathrm{KZ}}_{ij}\) is the KZ-normalised Casimir acting
on the \(i\)-th and \(j\)-th tensor factors of \(V^{\otimes n}\), and
\(dz_{ij}=dz_i-dz_j\).
The connection is flat: its curvature vanishes by the
infinitesimal braid relation
\begin{equation}\label{eq:ibr}
  [\Omega_{ij},\, \Omega_{ik} + \Omega_{jk}] \;=\; 0
  \qquad \text{for distinct } i, j, k.
\end{equation}

The monodromy of $\nabla_{\mathrm{KZ}}$ at $n = 2$ on the
configuration space $\Conf_2(\CC)$ is the single operator
\begin{equation}\label{eq:kz-monodromy-n2}
  \mathrm{Mon}(\nabla_{\mathrm{KZ}})
  \;=\;
  \exp\!\bigl(2\pi i\,\Omega^{\mathrm{KZ}}_\fg/(k+h^\vee)\bigr)
  \;\in\; \Aut(V \otimes V).
\end{equation}

\subsubsection{The Drinfeld associator and pentagon}

At $n = 3$, the monodromy of $\nabla_{\mathrm{KZ}}$ on
$\Conf_3(\PP^1 \setminus \{0,1,\infty\})$ produces the
Drinfeld associator
\begin{equation}\label{eq:drinfeld-associator}
  \Phi_{\mathrm{KZ}}
  \;=\;
  \Phi\!\bigl(\Omega^{\mathrm{KZ}}_{12}/(k+h^\vee),\,
  \Omega^{\mathrm{KZ}}_{23}/(k+h^\vee)\bigr)
  \;\in\; \exp\bigl(\hat{\mathfrak{t}}_3\bigr)
  \;\subset\; \End(V^{\otimes 3}),
\end{equation}
where $\hat{\mathfrak{t}}_3$ is the completed Drinfeld--Kohno Lie
algebra generated by $\Omega_{12}$ and $\Omega_{23}$ subject to
the infinitesimal braid relation~\eqref{eq:ibr}. The associator
$\Phi_{\mathrm{KZ}}$ satisfies the pentagon equation in $\gEone$:
\begin{equation}\label{eq:pentagon}
  \Phi_{1,2,34}\;
  \Phi_{12,3,4}
  \;=\;
  \Phi_{2,3,4}\;
  \Phi_{1,23,4}\;
  \Phi_{1,2,3},
\end{equation}
where $\Phi_{i,j,k}$ denotes the associator acting on the
$(i,j,k)$ tensor factors. This pentagon is the degree-$3$ MC
equation on the ordered side (the five faces of the Stasheff
associahedron $K_4$).

The antisymmetric linearised commutator
$[\Omega_{12}, \Omega_{23}]$ is killed by $\av_R$. The full
associator has the nonzero trivial-isotypic projection
\(\av_R(\Phi_{\mathrm{KZ}})=\mathfrak C(\cA)\) in general. The scalar
cubic shadow records this projection, while the ordered associator
contains the non-coinvariant data needed for the pentagon.

\subsubsection{Worked example: $\fg = \mathfrak{sl}_2$}
\label{subsubsec:sl2-example}

For $\fg = \mathfrak{sl}_2$ we have $d = 3$ and $h^\vee = 2$.
The modular characteristic is
\begin{equation}\label{eq:sl2-kappa}
  \kappa(\widehat{\mathfrak{sl}}_2,\, k)
  \;=\;
  \frac{3(k+2)}{4}.
\end{equation}
At $k = 0$: $\kappa = 3/2 = d/2$. At $k = -2$ (critical):
$\kappa = 0$.

The Yang $R$-matrix in the $2$-dimensional fundamental
representation $V = \CC^2$ is
\begin{equation}\label{eq:sl2-yang-r}
  R(u)
  \;=\;
  u\, I_4 + \hbar\, P
  \;\in\;
  \End(\CC^2 \otimes \CC^2),
  \qquad
  \hbar = \frac{1}{k + 2},
\end{equation}
where $P$ is the permutation matrix. In the standard basis
$\{e_1 \otimes e_1,\; e_1 \otimes e_2,\;
e_2 \otimes e_1,\; e_2 \otimes e_2\}$, this is the
$4 \times 4$ matrix
\begin{equation}\label{eq:sl2-r-matrix-4x4}
  R(u)
  \;=\;
  \begin{pmatrix}
    u + \hbar & 0 & 0 & 0 \\
    0 & u & \hbar & 0 \\
    0 & \hbar & u & 0 \\
    0 & 0 & 0 & u + \hbar
  \end{pmatrix}.
\end{equation}
The Yang--Baxter equation
$R_{12}(u-v)\, R_{13}(u)\, R_{23}(v)
= R_{23}(v)\, R_{13}(u)\, R_{12}(u-v)$
is verified by direct computation on $(\CC^2)^{\otimes 3}$.
The unitarity relation is $R(u)\, R(-u) = (u^2 - \hbar^2)\, I_4$.

The RTT relation
$R_{12}(u-v)\, T_1(u)\, T_2(v) = T_2(v)\, T_1(u)\, R_{12}(u-v)$
at level~$1$ gives the commutation relations of the Yangian
$Y_\hbar(\mathfrak{sl}_2)$:
\begin{equation}\label{eq:sl2-rtt-relations}
  [t_{ij}^{(1)},\, t_{kl}^{(1)}]
  \;=\;
  \hbar\bigl(\delta_{kj}\, t_{il}^{(1)}
  - \delta_{il}\, t_{kj}^{(1)}\bigr),
\end{equation}
which is $\mathfrak{sl}_2$ at $\hbar \neq 0$ (recovering the
ordinary Lie bracket $[e_{ij}, e_{kl}]
= \delta_{kj} e_{il} - \delta_{il} e_{kj}$), and becomes
abelian at $\hbar = 0$.

The KZ monodromy on $\Conf_2(\CC)$ decomposes $V \otimes V
= \CC^2 \otimes \CC^2$ into the singlet $V_0$ (dimension~$1$)
and the triplet $V_1$ (dimension~$3$), with eigenvalues
\begin{equation}\label{eq:sl2-monodromy-eigenvalues}
  \exp\!\bigl(2\pi i\,\Omega^{\mathrm{KZ}}_{\mathfrak{sl}_2}/(k+2)\bigr)
  \;=\;
  \begin{cases}
    q_{\mathrm{KL}}^{-3/2} & \text{on } V_0, \\[3pt]
    q_{\mathrm{KL}}^{1/2} & \text{on } V_1,
  \end{cases}
  \qquad q_{\mathrm{KL}}=e^{\pi i/(k+2)}.
\end{equation}
At integer level $k \in \ZZ_{\geq 1}$, these are roots of unity,
and the representation category of $U_q(\mathfrak{sl}_2)$
at $q = e^{\pi i/(k+2)}$ is semisimple with $k+1$ integrable
modules $V_0, \ldots, V_k$.

% ================================================================
% SECTION 9
% ================================================================
\section{The chiral quantum group equivalence}
\label{sec:chiral-qg}

On the Koszul locus, three structures on an $\Eone$-chiral algebra
determine each other: the $R$-matrix, the $\Ainf$ structure, and the
spectral coproduct. This section records the equivalence theorem,
extends it to the $\mathfrak{gl}_N$ tower for all $N \geq 1$,
exhibits the universal Miura coefficient, and recovers the
Verlinde formula from ordered chiral homology.

\subsection{The equivalence theorem}

\begin{theorem}[Chiral quantum group equivalence]
\label{thm:chiral-qg-equiv-standalone}
On the Koszul locus, the following three structures on an
$\Eone$-chiral algebra $\cA$ determine each other, up to the
choice of a Drinfeld associator $\Phi$:
\begin{enumerate}[label=\textup{(\Roman*)}]
\item \textup{($R$-matrix.)}
  A solution $r(z) \in \End_\cA(2) \otimes \cO(*\Delta)$ of
  the classical Yang--Baxter equation (the degree-$2$ MC
  component of $\Theta^{\Eone}_\cA$).

\item \textup{($\Ainf$-chiral structure.)}
  The full MC element $\Theta^{\Eone}_\cA$ equips the
  ordered bar complex $\Barord(\cA)$ with an
  $\Ainf$-chiral coalgebra structure: the higher products
  $m_k^{\mathrm{ch}}$ for $k \geq 2$ are extracted from the
  OPE on the real locus
  $K_k \hookrightarrow \overline{\FM}_k^{\mathrm{ord}}(\CC)$.

\item \textup{(Spectral coproduct.)}
  A spectral Drinfeld coproduct
  $\Delta_z\colon Y(\cA) \to Y(\cA) \otimes Y(\cA)$
  on the ordered Koszul dual Yangian, satisfying the
  \emph{associator-twisted} spectral coassociativity identity
  \begin{equation}\label{eq:spectral-coassoc}
    \Phi_{12,3} \circ (\Delta_{z_1} \otimes \id) \circ \Delta_{z_1+z_2}
    \;=\;
    (\id \otimes \Delta_{z_2}) \circ \Delta_{z_1} \circ \Phi_{1,23},
  \end{equation}
  where
  $\Phi = \Phi_{\mathrm{KZ}} \in U(\hat{\mathfrak g})^{\otimes 3}[[\hbar]]$
  is the Kontsevich--Drinfeld KZ associator, and
  $\Phi_{12,3}$, $\Phi_{1,23}$ denote its two codegeneracies
  acting on the triple tensor product.
\end{enumerate}

\begin{remark}[Pentagon cocycle and hexagon]
\label{rem:kz-associator-pentagon}
The associator $\Phi_{\mathrm{KZ}}$ appearing in
\eqref{eq:spectral-coassoc} is a \emph{pentagon cocycle}:
it satisfies Drinfeld's pentagon axiom
\[
  \Phi_{12,3,4} \cdot \Phi_{1,2,34}
  \;=\;
  (\id \otimes \Phi_{2,3,4}) \cdot \Phi_{1,23,4}
  \cdot (\Phi_{1,2,3} \otimes \id)
\]
in $U(\hat{\mathfrak g})^{\otimes 4}[[\hbar]]$, together with the
two hexagon axioms relating $\Phi$ to the $R$-matrix
$R(z) = e^{\hbar r(z)/2}$ of clause~(I)
\cite{Drinfeld91,Kontsevich99}. Equivalently,
$\Phi_{\mathrm{KZ}}$ is the holonomy of the Knizhnik--Zamolodchikov
connection on $\overline{\FM}_3^{\mathrm{ord}}(\CC)$ between the
limiting tangential basepoints $(z_1 \to z_2) \to z_3$ and
$z_1 \to (z_2 \to z_3)$; strictness of
\eqref{eq:spectral-coassoc} would correspond to $\Phi = 1$, which
fails already at order $\hbar^2$ by $\zeta(2)$. Ordered chiral
homology therefore carries an intrinsic Grothendieck--Teichm\"uller
action (cf.\ the GRT-parametrized face enumeration, Vol II Seven
Faces Theorem).
\end{remark}
The equivalences form a triangle:
\begin{equation}\label{eq:equiv-triangle-standalone}
\begin{tikzcd}[column sep=3em, row sep=3em]
  \textup{(I)} \ar[rr, shift left=1, "\alpha"]
  \ar[dr, shift left=1, "\gamma^{-1}"']
  && \textup{(II)} \ar[ll, shift left=1, "\alpha^{-1}"]
  \ar[dl, shift left=1, "\beta"] \\
  & \textup{(III)}
  \ar[ul, shift left=1, "\gamma"]
  \ar[ur, shift left=1, "\beta^{-1}"']
\end{tikzcd}
\end{equation}
The direction $\textup{(I)} \to \textup{(II)}$
($\alpha$: $R$-matrix determines $\Ainf$) is by restriction of
the OPE to the real Stasheff locus $K_n \subset
\overline{\FM}_n^{\mathrm{ord}}(\CC)$.
The direction $\textup{(II)} \to \textup{(III)}$
($\beta$: $\Ainf$ determines coproduct) is by cobar dualisation
of the bar coalgebra.
The direction $\textup{(III)} \to \textup{(I)}$
($\gamma$: coproduct determines $R$-matrix) is by extracting
the degree-$2$ component of $\Delta_z$ at the classical level.
\end{theorem}

\begin{proof}
$\alpha$: the Stasheff associahedron $K_n$ is isomorphic to
the real Fulton--MacPherson compactification
$\overline{\FM}_n^{\mathrm{ord}}(\RR)$
(Sinha~\cite{Sinha04}). The chiral $\Ainf$-maps
$m_k^{\mathrm{ch}}$ are defined by integrating the OPE against
the canonical volume form on $K_k$:
\[
  m_k^{\mathrm{ch}}(a_1, \ldots, a_k)
  \;=\;
  \int_{K_k} Y(a_1, z_1) \cdots Y(a_k, z_k)\, \omega_k.
\]
The Stasheff identity follows from Stokes' theorem on
$\partial K_{n+1}$.

$\beta$: the cobar construction
$\Omega^{\mathrm{ch}}(\Barord(\cA))$ is an $\Ainf$-algebra
whose universal property provides $\Delta_z$ by adjunction.
The spectral parameter $z$ arises from the meromorphic
dependence of the bar differential on the collision coordinate.

$\gamma$: the coproduct $\Delta_z$ at leading order in $z^{-1}$
gives $\Delta_z(T(u)) = T(u) \otimes T(u-z)$, and the
coefficient of $z^{-1}$ in the cross-term recovers $r(z)$.

The equivalence holds up to the choice of Drinfeld associator
$\Phi$: the cohomological derived centre is
associator-independent (proved for $\mathfrak{sl}_2$: the $\GRT_1$
action shifts the degree-$3$ obstruction class by coboundaries),
while the cochain-level derived centre is
associator-dependent. Bar-side invariants ($\kappa$, shadow tower)
are associator-free; cobar-side and coproduct-side depend
on~$\Phi$. See~\cite[Chapter~9]{Lor26} for the full proof.
\end{proof}

\begin{remark}[Verification scope]
\label{rem:verification-scope}
The equivalence is proved abstractly on the Koszul locus.
Concrete verification exists for $\mathfrak{sl}_2$ Yangian and
affine Kac--Moody. The elliptic case (KZB connection at
genus~$1$) is partial: the KZB local system is constructed, but the
full quantum group equivalence at genus~$1$ is open. The toroidal
case (genus~$2$ and higher) is absent.
\end{remark}

\subsection{The $\mathfrak{gl}_N$ tower}

The equivalence theorem extends to a uniform tower indexed by
$N \geq 1$, through the principal $\cW$-algebras
$\cW_N[\Psi]$ at level~$\Psi$.

\begin{theorem}[$\mathfrak{gl}_N$ chiral quantum group tower]
\label{thm:glN-chiral-qg-standalone}
For $N=1$ the datum is Heisenberg, and for $N=2$ it is the
Virasoro/$\mathfrak{sl}_2$ case. For $N\geq3$ and generic
level~$\Psi$, the chiral quantum-group datum for
$\cW_N[\Psi]$ is conditional on the Jindal--Kaubrys--Latyntsev vertex
bialgebra theorem~\cite{JKL26}. On the corresponding locus:
\begin{enumerate}[label=\textup{(\roman*)}]
\item \textup{(Transfer matrix.)}
  The $N \times N$ transfer matrix
  \begin{equation}\label{eq:glN-transfer-standalone}
    T(u)
    \;=\; I_N + \sum_{r \geq 1} T^{(r)}\, u^{-r},
    \qquad
    (T^{(r)})_{ij} = t_{ij}^{(r)},
    \quad 1 \leq i,j \leq N,
  \end{equation}
  is obtained from the quantum Miura factorisation
  $T(u) = (u + \Lambda_1) \cdots (u + \Lambda_N)$,
  where $\Lambda_i$ are the Miura fields.

\item \textup{(Yang's rational $R$-matrix.)}
  The $R$-matrix for $Y(\widehat{\mathfrak{gl}}_N)$
  in the fundamental representation is
  \begin{equation}\label{eq:yang-r-matrix-standalone}
    R(u)
    \;=\; u\, I_{N^2} + \Psi\, P
    \;\in\; \End(\CC^N \otimes \CC^N),
  \end{equation}
  where $P(e_i \otimes e_j) = e_j \otimes e_i$ is the
  permutation and $\Psi$ is the level prefix. The classical
  $r$-matrix is
  \begin{equation}\label{eq:classical-r-glN}
    r(z)
    \;=\;
    \Psi\, P/z,
  \end{equation}
  with $r|_{\Psi = 0} = 0$.

\item \textup{(Spectral Drinfeld coproduct.)}
  The coproduct $\Delta_z$ acts on the generators by
  \begin{equation}\label{eq:glN-coproduct}
    \Delta_z\bigl(T(u)\bigr)
    \;=\;
    T(u) \otimes T(u - z),
  \end{equation}
  and on the spin-$s$ primary $W_s$ ($s \geq 2$) by
  the cross-term formula of
  Theorem~\textup{\ref{thm:miura-standalone}} below.
\end{enumerate}
At $N = 1$: $\cW_1 = \cH_\Psi$ (Heisenberg at level $\Psi$),
the transfer matrix is scalar, $R(u) = u + \Psi$, and the
coproduct is an algebra homomorphism (no cross-terms).
At $N = 2$: $\cW_2 = \Vir_c$ with
$c = 1 - 6(\Psi - 1)^2/\Psi$, and the RTT relation at level~$1$
recovers $\mathfrak{sl}_2$. The Yang--Baxter equation
$R_{12}(u-v)\, R_{13}(u)\, R_{23}(v)
= R_{23}(v)\, R_{13}(u)\, R_{12}(u-v)$
holds for all $N$ and all $\Psi$.
\end{theorem}

\subsection{The universal Miura coefficient}

\begin{theorem}[Universal Miura cross-term coefficient]
\label{thm:miura-standalone}
Let $W_s$ denote the spin-$s$ primary generator of
$\cW_{1+\infty}[\Psi]$ ($W_1 = J$, $W_2 = T$, $W_3 = W$,
$\ldots$). For each $s \geq 2$, the spectral coproduct satisfies
\begin{equation}\label{eq:miura-universal}
  \Delta_z(W_s)
  \;=\;
  W_s \otimes 1 + 1 \otimes W_s
  + \frac{\Psi - 1}{\Psi}
  \bigl(J \otimes W_{s-1} + W_{s-1} \otimes J\bigr)
  + (\textup{lower-spin composites})
  + (\textup{spectral tail in~$z$}).
\end{equation}
The coefficient $(\Psi - 1)/\Psi$ of the primary cross-term
$J \otimes W_{s-1} + W_{s-1} \otimes J$ is
\emph{independent of the spin~$s$}.
\end{theorem}

\begin{proof}
The argument proceeds in three steps.

\emph{Step~1: the single-$J$ Miura coefficient is
$1/\Psi$.}
The quantum Miura factorisation writes
$\psi_s = e_s(\Lambda_1, \Lambda_2, \ldots)$, the $s$-th
elementary symmetric function of the Miura fields. Each
$\Lambda_i$ decomposes as $\Lambda_i = \sigma_i + J_i/\Psi$.
Expanding $e_s$ and collecting the sector with exactly one
$J$-insertion gives a coefficient of $1/\Psi$ on the
${:}J \cdot W_{s-1}{:}$ composite at every spin~$s$.

\emph{Step~2: the Drinfeld shift adds $+1$.}
The Drinfeld coproduct
$\Delta_z(\Lambda_i) = \Lambda_i \otimes 1 + 1 \otimes \Lambda_i$
is an algebra homomorphism on the $\psi$-generators. On the
$W$-generators, the Miura nonlinearity
$W_s = \psi_s - (\text{nonlinear correction})$
shifts the coefficient from $1/\Psi$ (the Miura coefficient of
${:}J \cdot W_{s-1}{:}$ in $\psi_s$) to
$(1/\Psi) + 1 - 1 = 1/\Psi$. The correction arises from
$\Delta_z(W_s) - \Delta_z(\psi_s)
= -\Delta_z(\text{nonlinear correction})$,
and the nonlinear correction at the single-$J$ level contributes
$-(1/\Psi) + (\Psi - 1)/\Psi = (1/\Psi)(\Psi - 1 - 1 + 1)$.
The net coefficient is $(\Psi - 1)/\Psi$.

\emph{Step~3: lower sectors do not contribute.}
The sectors with two or more $J$-insertions contribute to
composite terms $J \otimes J \otimes W_{s-2}$, etc., which
are the ``lower-spin composites'' in~\eqref{eq:miura-universal}.
These do not affect the primary cross-term coefficient.

The universality (independence of~$s$) is structural: it
follows from the fact that the Miura factorisation is a ring
homomorphism, so the single-$J$ coefficient is determined
by the linear algebra of $e_s$ at all~$s$ simultaneously.
Numerical verification: the coefficient $(\Psi-1)/\Psi$ has been
checked at spins $s = 2, 3, 4, 5, 6$ by direct computation.
\end{proof}

\begin{remark}[Coincidental agreement at $\Psi = 2$]
\label{rem:psi-2-coincidence}
At $\Psi = 2$, the coefficient $(\Psi - 1)/\Psi = 1/2 = 1/\Psi$:
the Miura coefficient and the coproduct coefficient agree
\emph{accidentally}. This coincidence masks the structural
distinction between $1/\Psi$ (the Miura contribution from
$e_s(\Lambda_i)$) and $(\Psi - 1)/\Psi$ (the coproduct
coefficient on $W$-generators). Verification at $\Psi = 3$
breaks the degeneracy: $1/3 \neq 2/3$.
\end{remark}

\subsection{Verlinde recovery}

At integer level $k \geq 1$, the ordered chiral homology of
$V_k(\mathfrak{sl}_2)$ on a genus-$g$ curve recovers the
Verlinde formula.

\begin{proposition}[Verlinde formula from ordered chiral homology]
\label{prop:verlinde-standalone}
Let $k \geq 1$ be a positive integer and
$S_{jl} = \sqrt{2/(k+2)}\,\sin(\pi(j+1)(l+1)/(k+2))$
the modular $S$-matrix for $\widehat{\mathfrak{sl}}_2$ at
level~$k$.
\begin{enumerate}[label=\textup{(\roman*)}]
\item \textup{(Genus~$0$.)} The ordered chiral homology
  recovers the unique vacuum:
  \begin{equation}\label{eq:verlinde-g0-standalone}
    Z_0(k)
    \;=\;
    \sum_{j=0}^{k} S_{0j}^2 \;=\; 1.
  \end{equation}

\item \textup{(Genus~$1$.)} The ordered chiral homology
  recovers the number of integrable representations:
  \begin{equation}\label{eq:verlinde-g1-standalone}
    Z_1(k) \;=\; \sum_{j=0}^{k} S_{0j}^0 \;=\; k+1.
  \end{equation}

\item \textup{(General genus.)} At genus $g \geq 2$:
  \begin{equation}\label{eq:verlinde-general-standalone}
    Z_g(k)
    \;:=\;
    \sum_{j=0}^{k} S_{0j}^{2-2g}.
  \end{equation}
  This is an integer for all $g \geq 0$, $k \geq 1$.
  At $k = 1$: $S_{00} = S_{01} = 1/\sqrt{2}$, so
  $Z_g(1) = 2 \cdot (1/\sqrt{2})^{2-2g} = 2^g$.

\item \textup{(Handle attachment.)}
  The ordered bar factorisation under non-separating sewing
  recovers the TQFT handle-attachment formula: each
  isospin channel $j$ contributes $H_j = S_{0j}^{-2}$, so
  \begin{equation}\label{eq:handle-standalone}
    Z_{g+1}(k)
    \;=\;
    \sum_{j=0}^{k} S_{0j}^{-2}\, S_{0j}^{2-2g}
    \;=\;
    \sum_{j=0}^{k} S_{0j}^{-2g}.
  \end{equation}
\end{enumerate}
\end{proposition}

\begin{proof}
The ordered chiral chain complex of $V_k(\mathfrak{sl}_2)$ on
$\Sigma_g$ is computed using the KZ local system of
Section~\ref{sec:heis-km}. At integer level, the representation
category of $U_q(\mathfrak{sl}_2)$ at $q = e^{\pi i/(k+2)}$ is
semisimple with $k + 1$ integrable modules $V_0, \ldots, V_k$.
The KZ local system has finite monodromy, and the ordered bar
complex factorises over the handle-sewing boundary strata
of $\overline\cM_g$.

At genus~$0$: unitarity of the modular $S$-matrix gives
$\sum_l S_{0l}\,\overline{S_{0l}} = 1$; for the self-dual
$\widehat{\mathfrak{sl}}_2$ $S$-matrix this reduces to
$\sum_l S_{0l}^{\,2} = 1$, hence $Z_0 = 1$.
(The stronger statement $\sum_l S_{jl}\,\overline{S_{il}} = \delta_{ij}$
holds for every $i,j$; we use only the $i=j=0$ diagonal.)
At genus~$1$: $S_{0j}^0 = 1$ for all $j$, so $Z_1 = k + 1$,
matching the count of integrable representations via
the Zhu algebra $A(V_k(\mathfrak{sl}_2))
\cong U(\mathfrak{sl}_2)/(e^{k+1}, f^{k+1})$.
At genus $g \geq 2$: each isospin channel $j$ contributes
$S_{0j}^{2 - 2g}$ by iterating the handle-attachment formula
$g - 1$ times from $Z_1 = k + 1$, verifying factorisation of
the ordered chiral homology under the separating and
non-separating sewing maps.
\end{proof}

\subsection{The three tiers}

The three tiers of the $\Eone$/$\Einf$ hierarchy refine
the equivalence theorem.

\begin{enumerate}[label=\textup{(Tier \alph*)}]
\item \emph{Pole-free ($\Einf$-chiral, class~$G$).}
  The $R$-matrix is the Koszul-signed permutation
  $S(z) = (-1)^{|a||b|}\id$. The Yangian is trivial. The
  coproduct is constant ($z$-independent). The formality bridge
  (Theorem~\ref{thm:formality-bridge-standalone}) holds: ordered
  and symmetric give identical answers. The averaging map loses
  nothing.

\item \emph{Vertex-algebra ($\Einf$-chiral, class~$L$/$C$/$M$).}
  The $R$-matrix is derived from the OPE and is non-trivial
  ($r(z) \neq 0$), but the $\cD$-module still descends to
  $X^{(n)}$ because the algebra is $\Einf$. The Yangian
  $Y_\hbar(\fg)$ is the $\Eone$ Koszul dual. The formality
  bridge holds at the cohomological level; the chain-level
  $R$-matrix is lost under $\av$.

\item \emph{Genuinely $\Eone$-chiral (Yangians, EK quantum
  vertex algebras).}
  The $R$-matrix is an independent input, not derived from any
  $\Einf$-chiral OPE. The $\cD$-module does not descend to
  $X^{(n)}$ (Theorem~\ref{thm:formality-failure-standalone}).
  The symmetric bar does not exist. The ordered bar is the only
  available construction, and the equivalence is the chiral
  analogue of the Etingof--Kazhdan quantisation
  theorem~\cite{EK96, EK00}.
\end{enumerate}

\bigskip

\noindent
The ordered bar complex is the categorical logarithm of a chiral
algebra: it records the full perturbative content as ordered
collision data, and the bar-cobar inversion recovers the algebra
from its logarithm. The symmetric bar is the $\Sigma_n$-invariant
trace of this logarithm: a scalar shadow that retains
$\kappa(\cA)$ but discards the $R$-matrix, the KZ associator, and
every braid-group representation that makes the $\Eone$ theory
strictly richer. The averaging map $\av_R$ is the sole information
loss in the passage from the $\Eone$ primitive to the
$E_\infty$ shadow. The five theorems of modular Koszul duality are
the invariants that survive this loss: they are projections of a
richer $\Eone$ structure --- the operadic circle, the chiral
quantum-group equivalence, and the chain-level computational
presentation of the closed colour of the chiral
Deligne--Tamarkin Swiss-cheese pair
$(Z^{\mathrm{der}}_{\mathrm{ch}}(\cA),\, \cA)$ via
$Z^{\mathrm{der}}_{\mathrm{ch}}(\cA) \simeq
\mathrm{Hom}_{\mathrm{ch}}(\Barord(\cA),\, \cA)$. The bulk is the
closed colour $Z^{\mathrm{der}}_{\mathrm{ch}}(\cA)$, not the bar
that resolves it. The holomorphic--topological theory in three
real dimensions is the universal-collar consequence of the
Swiss-cheese pair, supplied by
Theorem~\ref{thm:swiss-cheese-chiral-DT-standalone}, not by
iterating the bar functor on the chiral algebra alone.


% ================================================================
% REFERENCES
% ================================================================
\begin{thebibliography}{99}

\bibitem{Arnold69}
V.\,I.~Arnold,
\textit{The cohomology ring of the group of dyed braids},
Mat.~Zametki \textbf{5} (1969), 227--231.

\bibitem{BD04}
A.~Beilinson and V.~Drinfeld,
\textit{Chiral algebras},
American Mathematical Society Colloquium Publications, vol.~51,
AMS, Providence, RI, 2004.

\bibitem{CG17}
K.~Costello and O.~Gwilliam,
\textit{Factorization algebras in quantum field theory},
vols.~1--2, New Mathematical Monographs, Cambridge University
Press, 2017/2021.

\bibitem{Drinfeld85}
V.\,G.~Drinfeld,
\textit{Hopf algebras and the quantum Yang--Baxter equation},
Soviet Math.~Dokl.\ \textbf{32} (1985), 254--258.

\bibitem{Drinfeld91}
V.\,G.~Drinfeld,
\textit{On quasitriangular quasi-Hopf algebras and a group closely
connected with $\mathrm{Gal}(\overline{\mathbb{Q}}/\mathbb{Q})$},
Leningrad Math.~J. \textbf{2} (1991), 829--860.

\bibitem{EK96}
P.~Etingof and D.~Kazhdan,
\textit{Quantization of Lie bialgebras, I},
Selecta Math.\ (N.S.) \textbf{2} (1996), 1--41.

\bibitem{EK00}
P.~Etingof and D.~Kazhdan,
\textit{Quantization of Lie bialgebras, V: Quantum vertex
operator algebras},
Selecta Math.\ (N.S.) \textbf{6} (2000), 105--130.

\bibitem{FG12}
J.~Francis and D.~Gaitsgory,
\textit{Chiral Koszul duality},
Selecta Math.\ (N.S.) \textbf{18} (2012), 27--87.

\bibitem{GJ94}
E.~Getzler and J.\,D.\,S.~Jones,
\textit{Operads, homotopy algebra and iterated integrals for
double loop spaces},
preprint, 1994, arXiv:hep-th/9403055.

\bibitem{GK98}
E.~Getzler and M.\,M.~Kapranov,
\textit{Modular operads},
Compositio Math.\ \textbf{110} (1998), 65--126.

\bibitem{GLZ22}
J.~Gui, H.~Li, and J.~Zeng,
\textit{Quadratic duality for chiral algebras},
Adv.~Math.\ \textbf{406} (2022), 108496.

\bibitem{JKL26}
A.~Jindal, M.~Kaubrys, and E.~Latyntsev,
\textit{Vertex bialgebras and chiral quantum groups},
preprint, 2026.

\bibitem{Kohno85}
T.~Kohno,
\textit{S\'erie de Poincar\'e--Koszul associ\'ee aux groupes
de tresses pures},
Invent.~Math.\ \textbf{82} (1985), 57--75.

\bibitem{Kontsevich99}
M.~Kontsevich,
\textit{Operads and motives in deformation quantization},
Lett.~Math.~Phys.\ \textbf{48} (1999), 35--72.

\bibitem{Lor26}
R.~Lorgat,
\textit{Modular Koszul Duality, Volume~I}, monograph, 2026.

\bibitem{LV12}
J.-L.~Loday and B.~Vallette,
\textit{Algebraic operads},
Grundlehren der mathematischen Wissenschaften, vol.~346,
Springer, Heidelberg, 2012.

\bibitem{Molev07}
A.~Molev,
\textit{Yangians and classical Lie algebras},
Mathematical Surveys and Monographs, vol.~143,
AMS, Providence, RI, 2007.

\bibitem{MS24}
F.~Malikov and V.~Schechtman,
\textit{Homotopy chiral algebras},
preprint, 2024.

\bibitem{PR19}
T.~Proch\'azka and M.~Rap\v{c}\'ak,
\textit{Webs of $\cW$-algebras},
J.~High Energy Phys.\ \textbf{2018} (2018), no.~11, 109.

\bibitem{Sinha04}
D.\,P.~Sinha,
\textit{Manifold-theoretic compactifications of configuration
spaces},
Selecta Math.\ (N.S.) \textbf{10} (2004), 391--428.

\bibitem{SheltonYuzvinsky1997}
B.~Shelton and S.~Yuzvinsky,
\textit{Koszul algebras from graphs and hyperplane arrangements},
J.\ London Math.\ Soc.\ \textbf{56} (1997), 477--490.

\bibitem{Stasheff63}
J.\,D.~Stasheff,
\textit{Homotopy associativity of $H$-spaces, I},
Trans.~Amer.~Math.~Soc.\ \textbf{108} (1963), 275--292.

\bibitem{Vallette20}
B.~Vallette,
\textit{Homotopy theory of homotopy algebras},
Ann.~Inst.~Fourier \textbf{70} (2020), 683--738.

\bibitem{Verlinde88}
E.~Verlinde,
\textit{Fusion rules and modular transformations in 2D conformal
field theory},
Nuclear Phys.~B \textbf{300} (1988), 360--376.

\bibitem{Voronov99}
A.\,A.~Voronov,
\textit{The Swiss-cheese operad},
Contemp.\ Math.\ \textbf{239} (1999), 365--373;
arXiv:math/9807037.

\bibitem{KS00}
M.~Kontsevich and Y.~Soibelman,
\textit{Deformations of algebras over operads and the Deligne
conjecture},
in \emph{Conf\'erence Mosh\'e Flato 1999, Vol.~I}, Math.\ Phys.\
Stud.\ \textbf{21}, Kluwer, 2000, 255--307;
arXiv:math/0001151.

\end{thebibliography}

\end{document}
